% !TEX program = xelatex
% 从涡旋坍缩到形式化验证：OpenAI 纳维–斯托克斯方程有限时间爆破工作的数学思想与技术脉络
% 作者：Trae-Work、DeepSeek-V4.1-Flash
\documentclass[11pt,a4paper]{article}

\usepackage[UTF8,fontset=windows,heading=true]{ctex}
\usepackage{geometry}
\geometry{top=2.6cm,bottom=2.6cm,left=2.6cm,right=2.6cm}
\usepackage{amsmath,amssymb,amsthm,mathtools}
\usepackage[table]{xcolor}
\usepackage{booktabs}
\usepackage{array}
\usepackage{longtable}
\usepackage{enumitem}
\usepackage{caption}
\usepackage{float}
\usepackage[most]{tcolorbox}
\usepackage{titlesec}
\usepackage{fancyhdr}
\usepackage{fontspec}
\usepackage{xurl}
\usepackage{hyperref}

% ---------- 配色：藏青 + 天蓝 + 浅底 ----------
\definecolor{navy}{HTML}{17324F}
\definecolor{sky}{HTML}{2B7BB9}
\definecolor{skysoft}{HTML}{6FA8CF}
\definecolor{lightbg}{HTML}{F1F5F9}
\definecolor{theader}{HTML}{DCE9F5}
\definecolor{rulegray}{HTML}{B9C6D4}
\definecolor{inkgray}{HTML}{55606C}

% ---------- 中英文间距 ----------
\xeCJKsetup{CJKecglue={\hskip 0.22em plus 0.08em minus 0.05em}}

% ---------- 段落与换行 ----------
\linespread{1.32}
\setlength{\parskip}{0.45em}
\setlength{\emergencystretch}{3em}
\sloppy

% ---------- hyperref ----------
\hypersetup{
  colorlinks=true,
  linkcolor=navy,
  citecolor=navy,
  urlcolor=sky,
  bookmarksnumbered=true,
  pdfstartview=FitH,
  pdftitle={从涡旋坍缩到形式化验证：OpenAI 纳维-斯托克斯方程有限时间爆破工作的数学思想与技术脉络},
  pdfauthor={Trae-Work, DeepSeek-V4.1-Flash},
  pdfsubject={Navier-Stokes, Millennium Prize Problem, Formal Verification}
}

% ---------- 标题样式 ----------
\titleformat{\section}{\normalfont\Large\bfseries\color{navy}}{\thesection}{0.7em}{}
\titleformat{\subsection}{\normalfont\large\bfseries\color{navy}}{\thesubsection}{0.7em}{}
\titleformat{\subsubsection}{\normalfont\normalsize\bfseries\color{sky}}{\thesubsubsection}{0.7em}{}
\titlespacing*{\section}{0pt}{2.1ex plus 0.6ex minus 0.2ex}{1.1ex}
\titlespacing*{\subsection}{0pt}{1.6ex plus 0.5ex minus 0.2ex}{0.8ex}
\titlespacing*{\subsubsection}{0pt}{1.3ex plus 0.4ex minus 0.2ex}{0.6ex}

% ---------- 目录 ----------
\renewcommand{\contentsname}{目\hspace{0.5em}录}
\renewcommand{\refname}{参考文献}
\setcounter{tocdepth}{2}
\usepackage{tocloft}
\renewcommand{\cftsecfont}{\bfseries\color{navy}}
\renewcommand{\cftsecpagefont}{\bfseries\color{navy}}
\renewcommand{\cftsubsecfont}{\color{sky}}
\renewcommand{\cftsubsecpagefont}{\color{sky}}
\renewcommand{\cftsecleader}{\color{rulegray}\cftdotfill{\cftdotsep}}
\setlength{\cftbeforesecskip}{2pt}
\setlength{\cftbeforesubsecskip}{1pt}

% ---------- 定理类环境 ----------
\theoremstyle{plain}
\newtheorem{theorem}{定理}
\newtheorem{lemma}{引理}
\newtheorem{proposition}{命题}
\newtheorem{corollary}{推论}
\theoremstyle{definition}
\newtheorem{definition}{定义}
\theoremstyle{remark}
\newtheorem{remark}{注}
\renewcommand{\proofname}{\bfseries 证明}
\renewcommand{\qedsymbol}{$\blacksquare$}

% ---------- 卡片样式 ----------
\newtcolorbox{keybox}[1][]{
  colback=lightbg, colframe=navy, boxrule=0pt, leftrule=3pt,
  arc=1pt, left=9pt, right=9pt, top=7pt, bottom=7pt, breakable,
  fontupper=\small, #1}
\newtcolorbox{principlebox}[1][]{
  colback=white, colframe=sky, boxrule=0pt, leftrule=3pt,
  arc=1pt, left=9pt, right=9pt, top=7pt, bottom=7pt, breakable,
  fontupper=\small, #1}
\newtcolorbox{quotebox}[1][]{
  colback=lightbg, colframe=skysoft, boxrule=0pt, leftrule=3pt,
  arc=1pt, left=9pt, right=9pt, top=7pt, bottom=7pt, breakable,
  fontupper=\small\itshape, #1}

% ---------- 表格 ----------
\renewcommand{\arraystretch}{1.35}
\newcommand{\hdrow}{\rowcolor{theader}}
\captionsetup{font=small,labelfont={bf,color=navy},labelsep=quad}
\captionsetup[table]{skip=6pt}

% ---------- 页眉页脚 ----------
\pagestyle{fancy}
\fancyhf{}
\fancyhead[L]{\small\color{inkgray}OpenAI 纳维--斯托克斯工作综述}
\fancyhead[R]{\small\color{inkgray}\thepage}
\renewcommand{\headrulewidth}{0.4pt}
\renewcommand{\headrule}{\hbox to\headwidth{\color{rulegray}\leaders\hrule height \headrulewidth\hfill}}

% ---------- 引用命令 ----------
\newcommand{\arxiv}[1]{\href{https://arxiv.org/abs/#1}{arXiv:#1}}
\newcommand{\NS}{Navier--Stokes}

\title{\vspace{-1.2cm}\color{navy}\bfseries 从涡旋坍缩到形式化验证\\[0.3em]
\Large OpenAI 纳维--斯托克斯方程有限时间爆破工作的\\[0.2em]数学思想与技术脉络}
\author{Trae-Work、DeepSeek-V4.1-Flash}
\date{2026 年 10 月 13 日}

\begin{document}

\maketitle
\thispagestyle{empty}

\begin{abstract}
\noindent 2026 年 9 月 8 日，OpenAI 公布了一份由内部模型与约一万个并发智能体在 88 小时内产出的手稿，声称在光滑外力下构造出三维不可压缩 \NS{} 方程从静止出发、动能始终有界的有限时间爆破解，并给出 Lean 形式化。本文系统梳理这一工作的数学思想脉络：其技术血脉来自 Córdoba 与 Martínez-Zoroa 的“残差博弈”纲领，经 Alpöge--Buckmaster 在光滑外力与 Euler 方程上的推进而收束于 \NS{}。文章依次考察四层核心数学内容——各向异性自相似坍缩涡旋的标度、以零角平均振荡脉冲的平均动量通量补足外力、逐阶修正以保持外力 $C^{\infty}$ 光滑、以及在环面 $\mathbb{T}^{2}$ 上反演方向导数算子的 Fourier 乘子估计——并对每一步给出可直接核验的推导。随后转入数学界的验证与约束：Terence Tao 的方法溯源、Constantin--Ignatova--Vicol 关于“各向异性 Type II 界加轴对称核心加解析外力即正则”的刚性定理（它约束了构造的外力形态而非否证结论），以及 Bastounis--Circelli--Hansen 指出的自然语言证明与 Lean 形式化之间的落差，其中引理 8.6 的 $m+4$ 与 $m+5$ 差异被作为形式化翻译非语义忠实的样本。本文对该争议、署名与训练数据问题坚持严格中立并陈，最后给出经边界修正后的判断：该工作首先是一条技术路线的收敛证据，其次才是对形式化验证认识论的一次压力测试，而非千禧年问题的终局裁判。

\vspace{0.5em}
\noindent\textbf{关键词}\quad 纳维--斯托克斯方程；千禧年问题；有限时间爆破；自相似坍缩；Reynolds 应力；形式化验证；Lean；自动形式化
\end{abstract}

\vspace{0.6em}
\begin{keybox}
\textbf{阅读约定}\quad 本文所有关于手稿技术内容的陈述，均以 OpenAI 公开手稿、Lean 仓库与已公开的独立文献为依据，并尽量标注来源编号；凡属对证明正确性的独立复核，本文均未进行，相关判断一律引用公开文献并注明其作者立场。涉及争议事实的段落，均逐项标明来源、以归属语气分别叙述，不代任何一方作结论。
\end{keybox}

\vspace{0.5em}
{\small
\setlength{\parskip}{2pt}
\tableofcontents}

\clearpage

% ============ 第一部分：摘要后的第 1--3 章 ============

\section{引论：一道题的四种问法与一条被选中的路}

\subsection{官方提法与四个备选}

克雷数学研究所在 2000 年把 \NS{} 方程的存在性与光滑性问题列为七个千禧年问题之一。查尔斯·费弗曼（Charles Fefferman）撰写的官方提法\cite{clay}并没有给出单一命题，而是并列了四条互不依赖的路径，声明“证明其中任何一条即可”。其中两条是存在性陈述：三维空间中（备选 A）与周期盒 $\mathbb{T}^{3}$ 中（备选 B），从任意光滑、速降且无散度的初值出发，方程都有整体光滑解，且动能一致有界；另外两条是反例陈述：空间情形（备选 C）与周期情形（备选 D）允许求解者同时指定一个光滑初值和一个光滑外力，只要外力及其全部导数满足速降条件，使得不存在具有一致有界动能的整体光滑解。

\begin{keybox}
四条备选之间不是逻辑否定的关系。备选 C 允许在外力容许的条件下取反例，它不是备选 A 的否定：A 的否定应是“某个无外力初值没有整体光滑有界能量解”，而 C 的反例额外挂上了一项光滑外力。因此，构造出 C 只说明“在某些光滑推动下方程会失稳”，并不裁决“没有任何外力时方程是否永远光滑”。
\end{keybox}

官方文本对此的解释只有一句话，即这四条“给求解者合理的余地，同时保留问题的核心”。这句话本身就承认了：外力是问题的一部分，而不是外挂的技巧。

\subsection{2026 年 9 月发生了什么}

按 OpenAI 官方说明\cite{openai-blog}与其手稿\cite{openai-paper}，事件的时间轮廓大致如下：公司自 8 月 28 日起训练一个内部模型；9 月 1 日听闻“两个千禧年问题已被解决”的传闻，遂组织大规模智能体系统对若干公开难题发起搜索；先由约一百个智能体用约五十小时解决无外力三维 Euler 方程的爆破问题，随后把资源集中到 \NS{}；约一万个并发智能体在 88 小时后（9 月 5 日）给出 \NS{} 爆破解，Lean 形式化与验证再由 GPT-6 Astra 花费约十七小时完成。9 月 8 日，随官方说明一同发布的还有 166 页手稿《Finite time blowup for Navier--Stokes》与公开的 Lean 仓库\cite{openai-lean}；同日，纽约大学柯朗研究所的 Tristan Buckmaster 与 Anthropic 的 Levent Alpöge 也公开了他们在外力作用下的 IPM、Boussinesq 与三维 Euler 爆破结果，并对“抢发”提出异议。9 月 11 日，克雷研究所表态称该问题“看上去已被解决”，但强调评估与授信的流程“刻意不匆忙”\cite{clay-news}。

\begin{table}[H]
\centering
\small
\caption{2026 年 9--10 月关键节点（据公开来源整理）}
\begin{tabular}{@{}>{\raggedright\arraybackslash}p{2.6cm}>{\raggedright\arraybackslash}p{11.4cm}@{}}
\toprule
\hdrow \textbf{日期} & \textbf{事件} \\
\midrule
8/15--8/22 & Alpöge--Buckmaster 获得光滑外力下 IPM/Boussinesq/三维 Euler 的爆破结果，并完成 Lean 验证。 \\
9/1 & OpenAI 听闻传闻，启动多智能体战役。 \\
9/5--9/6 & 约一万并发智能体经 88 小时产出 \NS{} 爆破证明；Lean 形式化再耗约 17 小时。 \\
9/7 & Tao 发表博文，梳理 Alpöge--Buckmaster 的方法谱系，指向 Córdoba--Martínez-Zoroa。 \\
9/8 & OpenAI 发布手稿与 Lean 仓库；Buckmaster 同日发布“抢发”声明。 \\
9/11 & 克雷研究所表态“看上去已被解决”，评估“刻意不匆忙”。 \\
9/17 & Constantin--Ignatova--Vicol 上传刚性定理预印本；Gowers 撰文解释为何未签署菲尔兹奖得主公开信。 \\
10/6 & Bastounis--Circelli--Hansen 上传关于自动形式化非语义忠实的预印本，以引理 8.6 为例。 \\
10/10 & 媒体报道 OpenAI 已放出 722 篇 AI 生成数学手稿，学界反弹加剧。 \\
\bottomrule
\end{tabular}
\end{table}

\subsection{本文的范围、立场与方法}

本文是一篇基于公开材料的综述，而不是对该证明的独立复核。凡涉及“证明是否正确”的判断，本文一律转述公开文献并标明其作者的立场与保留意见；凡涉及技术内容的陈述，则以手稿原文、Lean 仓库与已公开的独立文献为准，并给出可核验的出处编号。这样做有两条理由。其一，166 页手稿加六十余万行 Lean 代码，其完全消化需要相当的时间，任何在数日内给出的“已确认/已推翻”结论都不诚实。其二，本文关心的问题不是“是非”，而是“这项工作的数学思想从哪里来、用了什么技术、这些技术和它们的验证处在什么位置”——这类问题可以依据公开材料稳妥地回答。

\begin{principlebox}
\textbf{本文主张的主线}\quad 这项工作的数学内容，可以压缩成一句话：先写下想要的奇异解，再让外力恰好等于残差；全部的困难都落在“让残差连同它的一切导数光滑地穿过奇点”上。围绕这一句话，才长出涡旋的标度、脉冲的通量、逐阶的修正与环面上的乘子估计。理解了这句话，才知道后面每一块技术各自在做什么。
\end{principlebox}

\subsection{记号约定}

全文用 $u(x,t)$ 表示速度场，$p$（或 $\pi$）表示压力，$f$ 表示外力，$\nu>0$ 为黏性系数。三组方程分别写作
\begin{align}
\text{（\NS{}）}\quad & \partial_{t}u+(u\cdot\nabla)u=-\nabla p+\nu\Delta u+f,\qquad \nabla\cdot u=0;\\
\text{（Euler）}\quad & \partial_{t}u+(u\cdot\nabla)u=-\nabla p+f,\qquad \nabla\cdot u=0.
\end{align}
柱坐标取 $x=(r\cos\theta,r\sin\theta,z)$，$u_r,u_\theta,u_z$ 为对应的三个分量。记奇点时间前的剩余时间为 $\tau=1-t$，径向与轴向特征长度分别记为 $\ell_r,\ell_z$。$\|u\|_{L^{\infty}}$ 指速度的 $L^{\infty}$ 范数，$\|u\|_{L^{2}}$ 指点动能的平方根。书写上，$\mathcal{P}u$ 表示速度关于 $z$ 轴的角向平均（轴对称部分），$w=u-\mathcal{P}u$ 为其非轴对称部分。

\subsection{为什么“允许外力”不是文字游戏}

对不熟悉该问题的人，“允许外力”听起来像是把难题偷偷改简单了。要消除这种疑虑，得把费弗曼的措辞读细。

在备选 A 与 B 里，外力恒为零，命题针对\textbf{一切}容许初值作全称量化，要求解永远光滑；在备选 C 与 D 里，外力可以非零，但命题是存在性的——只要能指出\textbf{一组}容许的初值与外力，使得不存在一致有界能量的整体光滑解，就算完成。看似后者更弱，实则不然。困难恰恰来自对\textbf{外力本身}的约束：它必须在空间上速降（对 C，条件是 $\big|\partial_{x}^{\alpha}\partial_{t}^{m}f(x,t)\big|\le C_{\alpha mK}(1+|x|+t)^{-K}$ 对一切 $\alpha,m,K$ 成立；对 D，只需在时间上速降），也就是光滑、无奇性、可以取紧支。于是，把奇点“塞进外力”这条捷径被堵死了：解得靠自己发散，外力却必须始终是个规矩的函数。

这一条约束，把问题从“找出一条发散的解”改造成“找出一条发散的解，同时让它留下的残差规矩”。残差越大越好造，残差规矩越难造。2026 年工作的全部技术重量，几乎都压在后半句上。理解这一点，是理解后文所有标度、脉冲与修正的钥匙。

\begin{keybox}
\textbf{一句话}\quad 外力不是外挂的作弊器，它是问题的一部分，而且是被约束得最严的那一项。解题者被允许挑一个外力，却不被允许挑一个坏的外力。
\end{keybox}

\section{前史：反例系谱与真正的困难所在}

要理解 2026 年的构造，先要知道此前“反例系谱”已经长成什么样。这一节的用意不是罗列年表，而是标出困难的位置：一百年来，人们已经能造出各种各样的病态解，却一直造不出一个“外力光滑的”病态解。

\subsection{Leray 弱解与能量不等式}

1934 年，让·勒雷（Jean Leray）\cite{leray}证明三维 \NS{} 方程存在整体有限能量的弱解，并满足能量不等式
\begin{equation}
\|u(t)\|_{L^{2}}^{2}+2\nu\int_{0}^{t}\|\nabla u(s)\|_{L^{2}}^{2}\,ds\le\|u_{0}\|_{L^{2}}^{2}.
\end{equation}
弱解的存在性不成问题，成问题的是光滑性与唯一性。勒雷本人的构造允许解在有限时间内失去正则性，此后“再启动”为另一个弱解。因此，从勒雷起，真正未知的就是：光滑初值出发的那条解，能否一直保持光滑。

\subsection{CKN 部分正则性：奇点集不可能太大}

1982 年，Caffarelli、Kohn 与 Nirenberg\cite{ckn}证明“适当弱解”的奇点集在抛物尺度下的一维 Hausdorff 测度为零。这是关于奇点几何的第一条强约束：即便奇点存在，它们也只能坐落在一个极小的集合上。这条结果给后来的构造者出了一道隐含的题——如果要造奇点，最好把它造得足够“干净”，例如时间上孤立、空间上单点，否则会与 CKN 的约束冲突。

\subsection{条件正则性与临界范数}

Escauriaza、Seregin 与 Šverák\cite{ess}证明：若无外力 Cauchy 问题的解在 $L_{t}^{\infty}L_{x}^{3}$ 中一致有界，则解光滑。$L^{3}$ 恰好是方程在抛物伸缩下的临界空间，这类“临界范数有界即正则”的结果划出了爆破解必须越过的门槛：真出现爆破，速度必须在临界范数上无界。它同时告诉后来者，人为制造的爆破应当从“临界范数发散而低于临界的量保持有界”入手。

\subsection{Tao 的平均化方程：在简化模型里实现爆破}

Terence Tao\cite{tao-averaged}构造了一个经过平均化的 \NS{} 型方程，它保留原方程的非线性能量抵消，却允许在有限时间内爆破。这条结果的意义是方法论上的：它说明“爆破”本身不是不可企及的，真正卡住人的是原方程中那些被平均化抹去的强非线性项。此后人们开始系统地把爆破解的构造问题，转译为“如何在保留关键结构的前提下安排一种可控的不稳定性”。

\subsection{凸积分与非唯一性}

2019 年，Buckmaster 与 Vicol\cite{bv-nonunique}用凸积分方法证明有限能量弱解的非唯一性；随后 Albritton、Brué 与 Colombo\cite{abc}用相似变量下的不稳定涡旋，构造出初值为零、外力相同却彼此不同的 Leray--Hopf 解。后者尤其接近 2026 年的构造：它同样在一个“看起来应当光滑”的设定里制造出病态行为。差别在于，它的外力属于 $L_{t}^{1}L_{x}^{2}$，在初始时刻奇异；而千禧年问题的备选 C 要求外力在空间与时间上都光滑、速降。把外力的正则性从“奇异”提升到“光滑”，正是此后数年的主攻方向。

\begin{keybox}
\textbf{前史给出的关键判断}\quad 到 2020 年代中期，制造“某种意义下的爆破解”已经不难；难的是让构造过程中出现的那个外力不跟着一起变坏。换句话说，困难的坐标已经从“方程内部”平移到了“外力”这一项上。2026 年的工作，正是冲着这一项去的。
\end{keybox}

\subsection{Onsager 判据与能量级串：临界指数为何反复出现}

在前史里还有一条不能略过的线索，它是理解后文“尺度”与“可和性”为何处处卡在临界值的背景。Onsager 在 1949 年提出：三维湍流中，若速度场的 Hölder 指数 $s>1/3$，则能量守恒；若 $s<1/3$，则可能出现能量在有限时间内耗尽的“反常耗散”。$1/3$ 这个数正是使能量通量在尺度间守恒的临界指数。此后的数学工作把 Onsager 的猜想严格化：临界指数以上的解保持能量守恒（常数与 Eyink 等人），其下则存在能量非守恒的弱解（含本节的凸积分构造）。

这条线索与本文的关系有两层。第一层是术语上的：后文会反复出现“临界范数”“可和临界指数”“按尺度的通量”，这些都是同一个尺度的不同面孔，$L^{3}$、$B^{1/3}_{3,\infty}$、二维可和指数 $-3$，彼此只差一个由维度决定的常数。第二层是思路上的：爆破构造普遍依赖“能量在小尺度上不完全被黏性抹平”这一图景；而二次面（脉冲的 Reynolds 应力、涡层的剪切）正是从小尺度向大尺度回输能量的通道。2026 年 \NS{} 构造里“脉冲向背景补给动量”的安排，与湍流中“小尺度向大尺度回能”的逆级串，结构上是同一件事的两种说法。

\begin{principlebox}
\textbf{一个反复出现的数}\quad $1/3$、$L^{3}$、$-3$、$B^{1/3}_{3,\infty}$ 是同一个尺度在四套语言里的写法。凡是围绕该问题的技术文本，遇到这些数都应当本能地想到“临界”——它们标记的是“恰好能守恒、恰好能求和”的分界，而不是随手选的常数。
\end{principlebox}

\subsection{凸积分非唯一性与千禧年爆破是两件事}

前史里另有一处容易混淆的地方，值得单独澄清：\textbf{凸积分（convex integration）给出的非唯一性，与千禧年问题所问的爆破，是两件不同的事。}

凸积分证明的是“弱解不唯一”：同一初值（甚至同一外力）可以对应多个不同的有限能量弱解，其中一些能量不守恒。2019 年 Buckmaster--Vicol 的结果即属此类。它的要害在于解的正则性极低——弱到不满足能量等式，正是这种“低正则”给了构造自由度。

而千禧年备选 C 要求的是另一回事：要找一条在奇点前 $C^{\infty}$ 光滑、动能一致有界的解，让它在有限时间自发失去 $L^{\infty}$ 有界性。这里没有“弱解”的余地，也没有“能量等式失效”的余地。换言之，凸积分的自由来自“放弃正则性”，而 \NS{} 爆破的困难恰在于“不能放弃正则性”。

把这两个目标分开之后，就能看清 2026 年构造的定位：它既没有走弱解的路线，也没有走“数值找近似解再证明稳定”的路线，而是走了一条介于两者之间的第三条路——完全在光滑范畴内构造一个残差控制得住的奇异流。它的技术难度，正来自“既要完全光滑、又要真的发散”这一对要求本身。

\subsection{为什么二维已解决、三维未决：涡旋拉伸}

前史里还有一个最基本的问题没有正面回答：为什么同一组方程，二维的整体正则性早已被证明，三维却至今悬置？答案只有四个字——涡旋拉伸。

在二维，涡量 $\omega$ 是一个标量，沿流线守恒：$\partial_{t}\omega+u\cdot\nabla\omega=\nu\Delta\omega$。既然涡量既不能被拉伸、也不能被放大，临界范数自然保持有界，正则性随之而来（Clay 的官方提法把这一条归功于 Ladyzhenskaya 等人）。在三维，涡量是一个向量，方程里多出 $\omega\cdot\nabla u$ 这一项——涡管可以被流场拉长。拉长使涡量截面变细、强度变大，这一放大机制可能压过黏性的扩散效应，从而使临界范数失去控制。于是三维的问题不在于方程有何不同，而在于多了一个“可以把涡量越拉越强”的自由度。

这一观察为后文提供了两条线索。其一是机械性的：2026 年构造正是主动利用涡旋拉伸——核心沿轴向拉长、半径收缩，角向速度随之增大，与 $\omega\cdot\nabla u$ 的放大效应同气连枝。其二是概念性的：任何三维爆破解都必须让某个临界范数发散，而二维构造无论怎么设计都无法突破这条限制；因此 2026 年的证明在二维不可能有对应版本，这不是技术不足，而是维度本身设下的壁垒。

\section{思想脉络（一）：Córdoba--Martínez-Zoroa 的“残差博弈”}

\subsection{先选解，后定力：残差博弈}

设想一般的演化方程 $\partial_{t}U+\mathcal{N}(U)=f$。对任意光滑的 $U$，只要定义 $f:=\partial_{t}U+\mathcal{N}(U)$，$U$ 就自动是“带外力方程”的解。于是“存在一个使其爆破的外力”这句话在字面上是平凡的。全部内容都落在对 $f$ 的正则性要求上：$f$ 必须光滑、紧支，并连同一切导数穿过奇点。Diego Córdoba 与 Luis Martínez-Zoroa 在 2024 年的工作\cite{cmz-ipm}里把这件事说成了一场博弈：允许残差的各个分项发散，但不允许它们的和发散。后来的所有构造都在这条博弈规则下进行。

这解释了为什么这类证明往往冗长：不是因为想法复杂，而是因为要同时把“位移”“速度”“压力”“黏性”几项都调到恰好抵消的尺度上，其记账量极大。

\subsection{迭代放大与高阶修正}

Córdoba--Martínez-Zoroa 在二维不可压多孔介质方程（IPM）中给出了第一个完整样本\cite{cmz-ipm}：以光滑、紧支的源项，构造出有限能量的光滑解，在有限时间内产生奇点。其骨架包含两个相互咬合的机制。

其一是逐尺度放大。先在低频上造一个解，再加入一个高频、小幅的扰动；扰动在其线性化的演化方程中经历不稳定增长，把解推向更奇异的方向；然后在这个更奇异的解之上再叠一层更细的扰动。如此层层叠加，层与层之间的时间间隔可以求和为一个有限的 $T_{*}$，从而在有限时间里塞进无穷多层。这样，梯度与涡量在 $t\uparrow T_{*}$ 时发散。

其二是随层号提高修正阶数。每一层的高频 ansatz 在截断包络之外必然留下振荡误差；要保证外力光滑，就必须把误差抵消到足够高的阶。Córdoba 与 Martínez-Zoroa 的做法是让修正阶数 $J_{q}$ 随层号 $q$ 一起趋于无穷，使外力级数在每一个导数意义上都收敛。这条“阶数随尺度同步上升”的原则，是整套构造能给出 $C^{\infty}$ 外力而非仅有有限阶光滑的关键。

\subsection{从粗糙外力到光滑外力}

Córdoba 与 Martínez-Zoroa 的后续工作把 IPM 的思路推广到带耗散的 \NS{} 型方程与外力下的 Euler 方程\cite{cmz-euler}，但外力的正则性起初只达到局部适定类，还不是备选 C 所要求的光滑速降。把外力正则性推到 $C_{c}^{\infty}$，是留给下一棒的问题。

\subsection{Alpöge--Buckmaster 的推进}

2026 年 8 月，Levent Alpöge 与 Tristan Buckmaster 沿着这条纲领前进，在光滑外力下先后拿下 IPM、二维 Boussinesq 与三维不可压 Euler 的爆破，并以 Lean 完成验证。Buckmaster 在声明\cite{buckmaster-statement}中明确把纲领的归属记在 Córdoba 与 Martínez-Zoroa 名下，称该攻击路线的“基本想法”来自二人，甚至表示 Martínez-Zoroa 应当因此获得菲尔兹奖。Tao 在 9 月 7 日的博文\cite{tao-blog-sep7}中独立地给出了同样的溯源，并把这条路线概括为：在低频解上反复叠加高频小幅修正，使解在临近爆破时越来越奇异，同时保持外力良态。

对本文而言，重要的是这条路线的“角色分工”在 Alpöge--Buckmaster 那里出现了一次转变。在 Boussinesq 中，被放大的波本身构成奇点：每一层被指数放大的温度梯度，成为下一层更细波的仿射背景。而在 \NS{} 中，奇点是一个先写下来的光滑自相似坍缩，被放大的波退居为工具，用它们的平均动量通量去补足坍缩所缺的那份外力。这一分工的差异，正是第 6 章要展开的核心。

\subsection{Tao 的抽象迭代框架}

Tao 在同一篇博文里把上述思想抽象成一个不依赖具体方程的三段式。设已有低频（外力作用下的）解 $U_{0}$，欲在其上加一个高频小幅修正 $U=U_{0}+u$，代价是外力上多出一个修正 $g$。若能把 $u$ 的振幅做得相对大、而 $g$ 的振幅做得相对小，且让修正的频率随迭代次数迅速升高，就可以指望这一过程收敛到一个在有限时间爆破的解 $U$，而外力保持光滑。

Tao 进一步指出，要让“小扰动变大”，基本手段是让 $u$ 所满足的线性化方程
\begin{equation}
\partial_{t}u+\mathcal{L}u=g
\end{equation}
具有可利用的不稳定性：$u$ 可以在早期指数地小，而在临近爆破时被放大到足以破坏光滑性，同时又不至于大到使非线性分析失控。在 Boussinesq 里，这个线性化系统甚至有显式解——在原点附近速度关于空间线性、温度呈高频平面波——从而化为一组显式的振幅 ODE，其不稳定性可以直接读出。第 4 章将把这类“精确波”计算完整推一遍，因为它既是 Boussinesq 的引擎，也是 \NS{} 中脉冲生长与通量生成的理论依据。

\begin{principlebox}
\textbf{脉络小结}\quad Córdoba--Martínez-Zoroa 提供了纲领（残差博弈＋逐尺度放大＋阶数随尺度上升）；Alpöge--Buckmaster 把纲领推进到光滑外力与 Euler，并改变了被放大扰动的角色；Tao 把它抽象为“低频解＋可利用不稳定的高频修正”。OpenAI 的工作落在这一谱系的末端，其独有的贡献是把 \NS{} 的坍缩核心与“由离心不稳定性生长的脉冲所供给的 Reynolds 应力”精确地对接起来。
\end{principlebox}

\subsection{这条纲领的强项与短板}

把纲领摆到台面上，才能公正地评价 2026 年的推进。它的强项有三。其一，把“求解方程”转成了“设计抵消”，使问题从动力学变成可记账的代数，由此获得了极高的可操作性：一旦尺度对齐、阶数配好，剩下的就是逐项估计。其二，它对初值极不敏感，甚至可以从静止出发——因为解是被构造出来的，而初值只是一个待满足的边界条件。其三，它与外力正则性天然兼容：只要把修正阶数提到足够高，外力就能在任意多的导数意义下被压平。

它的短板同样清楚。第一，它给出的是存在性反例，而非对一般解的刻画：它说明“有些流会爆”，不说明“哪些流会爆”，更不触及无外力情形。第二，构造高度依赖人为设计，尺度、频率、修正阶数都需要精心选择，因而对“如果某一步的估计不够好会怎样”缺乏鲁棒性说明。第三，也是最容易被忽略的一点：这类构造的“可读性代价”很高——为了把外力压光滑，技术细节层层叠加，最终产物往往长达上百页，其核心思想被淹没在估计之中。第 7 章的逐阶修正与第 10 章的理解之争，追根溯源都来自这第三条短板。

% ============ 第二部分：第 4--8 章（技术核心，含内联证明链） ============

\section{思想脉络（二）：把爆破化为配平问题的第一性原理}

\subsection{本质问题}

第 3 章已经把困难的位置指了出来：能不能造出爆破解，取决于能不能让残差 $f=\partial_{t}u+(u\cdot\nabla)u-\nu\Delta u+\nabla p$ 保持光滑。把这个判断再往前推一步，就得到整个技术方案的第一性原理。

若随便写一个在 $t\uparrow1$ 时发散的速度场，几乎必然得到一个同样发散的外力：一个刚性旋转、角速度按 $(1-t)^{-1/2}$ 增长的玩具涡旋，其压力取 $\frac12 a(t)^{2}(x^{2}+y^{2})$ 恰好抵消离心项，但残差仍然按 $(1-t)^{-3/2}$ 发散。可见“发明一个奇异涡旋”是容易的，难的是让残差里的各个发散分量恰好相消。于是工作真正的形状被确定下来：不是解一个方程，而是设计一场精确的抵消。

\begin{keybox}
\textbf{第一性原理}\quad \NS{} 爆破构造的困难不在动力学而在代数：要在速度、压力、对流、黏性四项各自的量级都趋于无穷时，安排它们的和连同一切导数保持有界。全部技巧都是在为这场抵消寻找足够的自由度。
\end{keybox}

\subsection{仿射背景上的精确波}

抵消的自由度从哪来？来自非线性项本身的“不自作用”性质。考虑一段仿射背景 $u_{\text{old}}(x,t)=D(t)x$，其中 $\operatorname{tr}D=0$（保证无散），再叠加一列平面波
\begin{equation}
w=a(t)\,\phi\big(\xi(t)\cdot x\big),\qquad a\perp\xi,
\end{equation}
其中 $a\perp\xi$ 使 $\nabla\cdot w=(a\cdot\xi)\phi'=0$。此时波的对流项
\begin{equation}
(w\cdot\nabla)w=a(a\cdot\xi)\,\phi\,\phi'=0
\end{equation}
因为 $a\cdot\xi=0$ 而恒为零。这就是 Craik--Criminale 的精确波：波的二次自作用消失，波因而满足一个线性方程。非线性项里最麻烦的那一块被结构性地点掉了，剩下的动力学变成可控的线性系统。这类结构在流体中并不罕见——离心不稳定性、Rayleigh--Taylor 不稳定性都可以用精确波来承载。

\subsection{振幅方程：逐项推导}

为看清背景应变 $D$ 如何给波注入能量，把 $u=Dx+w$ 代入无外力 Euler 方程，并减去背景自身满足的方程。以下逐步列出存活项，这一步是后面一切定量估计的起点。

设背景满足 Euler：$\partial_{t}(Dx)+(Dx\cdot\nabla)(Dx)+\nabla p_{\text{old}}=0$。由于 $(Dx\cdot\nabla)(Dx)=D^{2}x$ 与 $\partial_{t}(Dx)=\dot Dx$ 都是关于 $x$ 线性的，背景压力必须也是 $x$ 的二次型，这与后面波的压力恰好相容。把 $u=Dx+w$ 代入并减去背景方程，剩余项为三部分：

其一，时间导数项 $\partial_{t}w=\dot a\,\phi+a\,\phi'\,(\dot\xi\cdot x)$。其二，波对流背景 $\big(w\cdot\nabla\big)(Dx)=Dw$，以及背景对流波 $\big(Dx\cdot\nabla\big)w=\phi'\,(Dx\cdot\xi)a$。其三，波的自身对流项已由 $a\cdot\xi=0$ 消去，只剩下压力项 $\nabla p_{w}=q(t)\,\xi\,\Phi(\xi\cdot x)$，其中 $\Phi'=\phi$，故压力梯度沿 $\xi$ 方向。

采用相位传输条件 $\dot\xi=-D^{\top}\xi$，则 $\partial_{t}w$ 中沿 $\phi'$ 的分量 $a\,(\dot\xi\cdot x)$ 与背景对流波项 $a\phi'\,(Dx\cdot\xi)$ 恰好相消——注意 $\dot\xi\cdot x=-(D^{\top}\xi)\cdot x=-\xi\cdot(Dx)$，两者之和为零。余下的 $\phi$ 系数项给出
\begin{equation}
\dot a=-Da-q\,\xi ,
\label{eq:coeffphi}
\end{equation}
其中 $q$ 是待定的压力标量。为使波的无散约束 $a\cdot\xi=0$ 在演化中保持，需要对它求时间导数并令其为零：$\dot a\cdot\xi+a\cdot\dot\xi=0$。代入 $\dot\xi=-D^{\top}\xi$，得 $a\cdot\dot\xi=-(Da)\cdot\xi$；再代入式 \eqref{eq:coeffphi}，得 $\dot a\cdot\xi=-(Da)\cdot\xi-q|\xi|^{2}$。两者相加为零，于是压力标量被唯一确定为
\begin{equation}
q=-2\,\frac{\xi\cdot Da}{|\xi|^{2}},
\end{equation}
而振幅满足
\begin{equation}
\dot a=-Da+2\,\frac{\xi(\xi\cdot Da)}{|\xi|^{2}}
=-\Big(I-\frac{2\,\xi\otimes\xi}{|\xi|^{2}}\Big)Da.
\label{eq:ampeq}
\end{equation}
矩阵 $I-\frac{2\xi\otimes\xi}{|\xi|^{2}}$ 是沿 $\xi^{\perp}$ 平面的反射：压力把 $-Da$ 中沿 $\xi$ 的分量抹去，又因 $\xi$ 自身在转动而反向补回。范数满足
\begin{equation}
\tfrac12\partial_{t}|\xi|^{2}=-\xi\cdot D\xi,\qquad \tfrac12\partial_{t}|a|^{2}=-a\cdot Da .
\end{equation}
这两式给出一个干净的判据：$\xi\cdot D\xi<0$ 使波数变大，但它本身并不保证振幅增长；振幅是否增长，本质上是背景 $D$ 的稳定性问题，可以从式 \eqref{eq:ampeq} 读出。这就是“精确波”这一想法的全部内容——把一个偏微分方程问题归约为一对常微分方程，再在这对常微分方程上寻找不稳定线。

\subsection{两种截然不同的用法}

这套归约在两篇工作里扮演了相反的角色，这个对照是理解 \NS{} 构造的关键。

在 Boussinesq 中，被放大的波本身就是奇点。取自由的 $D=0$、稳定的温度梯度 $G=-Ae_{2}$（$A>0$ 对应冷重流体压住暖轻流体的不稳定层结），写 $\zeta=r\,e(\varphi)$，则振幅方程给出二阶系统
\begin{equation}
\frac{d}{dt}\begin{pmatrix}\Theta\\ \Omega\end{pmatrix}
=\begin{pmatrix}0 & \dfrac{A\sin\varphi}{\lambda r}\\[2mm] \lambda r\sin\varphi & 0\end{pmatrix}
\begin{pmatrix}\Theta\\ \Omega\end{pmatrix},
\qquad\text{特征值}=\pm\sqrt{A}\,\sin\varphi .
\end{equation}
生长率 $\sqrt{A}\sin\varphi$ 与波数 $\lambda$ 无关，但原点处的梯度 $\nabla\vartheta(0,t)=\lambda\,\Theta\,\zeta$ 与 $\lambda$ 成正比。于是，一个振幅极小、频率极高的波可以在温度本身保持有界的同时，制造出无界的梯度——恰好是“$\|\theta\|_{\infty}$ 有界而 $\|\nabla\theta\|_{\infty}\to\infty$”所需的那种形状。代价是随之而来的剪切 $\Omega$，需要用共同旋转把它“转向”（旋转背景梯度与波矢，使实验室坐标下的 $\zeta_{1}$ 反号，把 $\Omega$ 拉回零），之后才能在其上叠加下一层更细的波。每一层因此经历“生长、过渡、转向、保持”四阶段，层间时间可求和。

在 \NS{} 中，奇点不再由波承担。构造者先写下一条光滑的自相似坍缩涡旋柱，让它自己去发散；被放大的波只是工具，任务是提供坍缩本身缺掉的那份外力。这个分工的落地方式就是下一节要讲的 Reynolds 应力。

\subsection{Boussinesq 级联的全景：从一层到无穷层}

上一小节给了单层的增长率，但“一层”不足以产生奇点，需要把无穷多层塞进有限时间。这一步的组织方式是：$\lambda_{q}=\lambda_{q-1}^{Q_{*}+q}$ 的超指数频率表、$\ell_{q}=\lambda_{q-1}^{-3}$ 的包络半径，以及 $\sigma_{q}\ge\sigma_{q-1}^{3}$ 的增长尺度。每层的生长时长约为 $\sigma_{q-1}^{-1}$，而 $\sigma_{q}$ 按立方增长，故各层时长求和收敛到一个有限的 $T_{*}$。于是无穷多层在 $T_{*}$ 之前落位，梯度在 $T_{*}$ 处发散。

要让这个级联真的收敛到“梯度发散而温度有界”，还需处理两件事。其一是增益不能被后续层抵消：所有梯度方向都落在以第一层方向为中心的一个锐角锥内，$\zeta_{q}$ 的方向被共同旋转“转向”后仍与原方向保持夹角下界，于是投影到 $-e_{1}$ 上的和仍有正的下界，$\|\nabla\theta(t)\|_{\infty}\ge cA_{n}(t)\to\infty$。其二是各层不能互相破坏：由于 $\ell_{q}\ll\ell_{q-1}$，更细的层被支撑在所有更粗层的仿射区域内，那里剧烈部分恰好是精确仿射的，高阶修正项为零，从而层与层之间的干扰被结构性排除。

这套“超指数频率＋锐角锥方向控制＋嵌套仿射区域”的三件套，是 Córdoba--Martínez-Zoroa 纲领的具体化身。理解了它，\NS{} 构造里那些看似琐碎的频率与阶数安排就有了出处——它们不是审美选择，而是收敛性的硬性要求。

\subsection{精确波为何是“引擎”}

把精确波称为引擎，不只是因为它提供增长，更因为它把“非线性爆破”这件难事，拆成了“线性稳定性”这件易事。爆破的时间、位置与速率，都可以从振幅 ODE 的特征值与相位方程直接读出：特征值 $\pm\sqrt{A}\sin\varphi$ 给出生长率，$\dot\xi=-D^{\top}\xi$ 给出波数随时间的走向，$\lambda\Theta\zeta$ 给出梯度的量级。换言之，一个原本要在偏微分方程层面讨论的问题，被压缩到几个标量 ODE 上。

这种压缩之所以可能，依赖两条结构性事实：波的极化与波矢正交（$a\perp\xi$，消去自作用），以及背景是仿射的（$\nabla(Dx)=D$ 是常数矩阵，使代数运算初等化）。一旦背景不再是仿射的，或波不再正交极化，压缩就会失效。于是这一工具的适用范围被划得很清楚：它服务于“局部仿射、局部平面波”的区域，其它区域要靠包络的截断与逐阶修正去处理。第 5、6 章分别把这两类区域用到核心与环域上。

\subsection{小结}

把这一章压缩成一句可记的话：精确波之所以有用，是因为它把“非线性”中最麻烦的自作用项结构性地消掉，从而把一类看起来无从下手的偏微分方程问题，变成一对可以手算稳定性的常微分方程；剩下的工作，就是决定让这对常微分方程的不稳定性去承担奇点，还是去承担外力。

\section{核心数学技术（一）：各向异性自相似坍缩涡旋}

\subsection{奇点的几何：一根收缩的细柱}

构造把奇点放在原点、时间归一化为 $t=1$。在柱坐标 $(r,\theta,z)$ 下，主导流 $u^{(0)}$ 是轴对称的，其分量只依赖 $(r,z,t)$。物理图像很直观：流体沿径向向内螺旋、同时沿轴向在 $z\approx0$ 分层面两侧向外排出；压力向轴心递减提供向心的主导力。角动量在向内输送中被黏性向外搬运，二者平衡决定了特征转速。

关键是自相似的标度。记 $\tau=1-t$，手稿取
\begin{equation}
\ell_{r}\asymp\tau^{1/2},\qquad \ell_{z}\asymp\tau^{1/2-h},\qquad 0<h<\tfrac1{100},
\label{eq:scales}
\end{equation}
其中 $\asymp$ 表示被与 $\tau$ 无关的正常数上下夹住。两个长度都趋于零，但
\begin{equation}
\frac{\ell_{r}}{\ell_{z}}\asymp\tau^{h}\to0,
\end{equation}
所以核心越来越细长，其体积为 $\ell_{r}^{2}\ell_{z}\asymp\tau^{3/2-h}$。

\subsection{速度标度与能量：局部无界与全局有界的相容}

与式 \eqref{eq:scales} 配套的速度量级为
\begin{equation}
|u_{\theta}^{(0)}|,\,|u_{z}^{(0)}|\asymp\tau^{-1/2-h},\qquad |u_{r}^{(0)}|=O(\tau^{-1/2}).
\end{equation}
现在直接估算核心的动能，看它为什么不发散：
\begin{equation}
E_{\text{core}}\;\sim\;\text{体积}\times(\text{特征速度})^{2}
\;\asymp\;\tau^{3/2-h}\cdot\tau^{-1-2h}
\;=\;\tau^{1/2-3h}\;\xrightarrow[\tau\to0]{}\;0 .
\label{eq:energy}
\end{equation}
因为 $h<1/100$，指数 $1/2-3h>0$ 仍为正，故核心动能趋于零——尽管速度在增长。这就是“局部速度无界、全局动能有界”之所以相容的算术根源：速度按 $\tau^{-1/2-h}$ 发散，但承载它的体积按 $\tau^{3/2-h}$ 收缩，二者相乘后能量反而变小。

\begin{keybox}
\textbf{一个可记的估算}\quad 无限速度之所以不违反有限能量，是因为它被压缩进了越来越小的体积。式 \eqref{eq:energy} 的指数 $1/2-3h$ 就是“能量不掉队”的全部秘密，$h$ 的上界 $1/100$ 正是为了保住这个不等式。
\end{keybox}

\subsection{两套 Reynolds 数}

用核心半径 $\ell_{r}$ 作长度、以角向与径向速度分别作速度，得
\begin{equation}
\mathrm{Re}_{\theta}=\frac{|u_{\theta}|\ell_{r}}{\nu}\asymp\tau^{-h}\to\infty,
\qquad
\mathrm{Re}_{r}=\frac{|u_{r}|\ell_{r}}{\nu}=O(1).
\end{equation}
角向 Reynolds 数无界增长，意味着流体在径向扩散时间内转了越来越多的圈；径向 Reynolds 数保持有界，意味着黏性始终与径向内流相抗衡。再看各项的变化率：
\begin{equation}
\frac{|u_{r}|}{\ell_{r}}=O(\tau^{-1}),\quad
\frac{|u_{z}|}{\ell_{z}}\asymp\tau^{-1},\quad
\frac{\nu}{\ell_{r}^{2}}\asymp\tau^{-1},
\qquad
\frac{\nu/\ell_{z}^{2}}{\nu/\ell_{r}^{2}}=\frac{\ell_{r}^{2}}{\ell_{z}^{2}}\asymp\tau^{2h}\to0 .
\end{equation}
前三个变化率都是 $\tau^{-1}$ 量级，说明径向扩散、径向对流与轴向对流同处主导平衡；轴向扩散相对径向扩散按 $\tau^{2h}$ 衰减而退居次席。这套尺度不是拍脑袋选的，而是“让主导项同阶、让次主导项可压”这一要求反解出来的。

\subsection{为什么必须打破反射对称}

若剖面关于 $z=0$ 严格对称，则角动量向中层面的轴向输运在 $z=0$ 处消失，且对称性迫使 $u_{z}(r,0,t)=0$，于是中层面附近既无转动放大、也无轴向速度的径向剪切可以补足。构造因此选取略不对称的轴向剖面：带一点向上的偏置、在 $z=0$ 处留一个非零的 $u_{z}$。这使“转动放大较弱的层”与“轴向剪切消失的层”在位置上分开，由 $u_{z}$ 的径向剪切在中层面附近补足所需的放大；稍远处，转动机制本身就够用。这一处看似琐碎的对称性破坏，是核心侧壁能把剪切喂给外侧脉冲的前提。

\subsection{与刚性定理假设的对应}

后文第 9 章要讨论的 Constantin--Ignatova--Vicol 刚性定理\cite{civ}，其假设正好抓住了本节的两个特征：其一，速度的角向平均 $\mathcal{P}u$ 满足各向异性的 Type II 界，即径向分量被 $(1-t)^{-1/2}$ 控制、轴向与角向分量被 $(1-t)^{-1/2-h}$ 控制，且每求一次 $r$ 导数多付 $\tau^{-1/2}$、每求一次 $z$ 导数只多付 $\tau^{-1/2+h}$；其二，速度在一个半径为 $c\sqrt{\tau}$ 的坍缩核心区严格轴对称。这两条与式 \eqref{eq:scales} 及核心构造一一对应。把这个对应先记住，第 9 章的结论才读得顺。

\subsection{核心内部的剖面方程与压力平衡}

标度回答的是“多快、多小”，剖面方程回答的是“什么形状”。在核心最内侧，主导动量平衡由一组关于 $(r,z)$ 的定常型方程给出；把它写成自相似变量后，径向与轴向的平衡各自成为一个常微分方程的边值问题（附录 B 专门处理轴附近的解析剖面与延拓）。物理上，压力向轴心递减，提供向内螺旋所需的向心力；轴向压力的梯度在中层面附近指向 $z=0$，与轴向输运和黏性共同维持轴向排出。由于不可压，向内汇集的流体必须沿轴向排出，否则会在轴心堆积——正是这条约束把“向内”、“旋转加速”、“轴向排出”三者锁在一起：径向内流携带角动量靠近轴心，转动因而加快；无扭矩时角动量 $ru_{\theta}$ 守恒，而实际核心中黏性把角动量向外搬走，特征转速由“向内搬运”与“黏性流失”的平衡决定。

这套剖面方程的存在性本身依赖附录 B 的解析性论证：只有在轴附近取到实解析的剖面，内层才能与外部流光滑对接。这也是“解析性”在本构造中第一次以正面身份出现——第 9 章刚性定理里，外力一旦在奇点附近解析就会逼出正则，而这里解析性只用于核心内部的剖面，两者作用方向相反，不可混为一谈。

\section{核心数学技术（二）：振荡脉冲与平均动量通量}

\subsection{环域里的“缺力”}

把自相似核心与一个速度随半径递减的光滑外部流对接，会在两者之间的环域留下一个不平衡。单看背景流，这个残差在 $t\uparrow1$ 时发散，无法充当所需的光滑外力。于是必须让流体自身的运动补上这份缺掉的动量输运——具体做法，是在环域布置一族空间振荡的脉冲。

\subsection{零角平均脉冲如何产生内部力}

脉冲的关键性质是：它们的角向平均为零，但其二次自作用不为零。对速度分量做角向平均，记 $\langle\cdot\rangle=\frac{1}{2\pi}\int_{0}^{2\pi}\cdot\,d\theta$，则对流项的平均包含
\begin{equation}
\big\langle (w\cdot\nabla)w\big\rangle
=\big(\langle w\rangle\cdot\nabla\big)\langle w\rangle
+\big\langle (w-\langle w\rangle)\cdot\nabla\,(w-\langle w\rangle)\big\rangle .
\end{equation}
第一项因 $\langle w\rangle=0$ 而消失，第二项一般不为零：两个平均为零的振荡分量，其乘积的平均不必为零。这一项正是 Reynolds 应力 $\langle w\otimes w\rangle$ 的散度，它充当作用在背景流上的内部力。选择脉冲的排布，使其提供的 Reynolds 应力恰好抵消背景残差的奇异部分，剩下的余项就能光滑地穿过 $t=1$。这一步是本项工作区别于前人的核心创新。

\begin{keybox}
\textbf{核心创新的一句话版}\quad 用平均为零的振荡脉冲，靠它们乘积的平均（Reynolds 应力）去补足坍缩涡旋在奇点处缺掉的那份动量输运；脉冲自己不发散，却替背景流承担了发散的外力。
\end{keybox}

\subsection{脉冲何以增长：离心不稳定性}

脉冲需要从“指数小”逐渐生长到足以承担通量。其生长机制是经典的离心不稳定性，可以追溯到 Leibovich--Stewartson 与 Billant--Gallaire 的离心不稳定判据\cite{leibovich,billant}。在环域中，角速度随半径向外递减，于是被径向推移的流体元处在不利于稳定的角动量分配中，扰动被进一步推离原位，形成指数增长。第 4 章的精确波归约在此再次出场：把脉冲看作叠加在缓慢变化的仿射背景上的精确波，其振幅 ODE 给出增长指数；脉冲的振荡性保证它对背景的“平均”效应正是我们要的 Reynolds 应力，而其局部起伏不破坏背景的光滑性。

\subsection{黏性阻尼与激活窗口}

与 Euler 方程不同，\NS{} 的脉冲要额外背负一项黏性阻尼。对正弦载波的精确波，方程中会多出一项 $-\nu|\xi|^{2}a$：频率越高，阻尼越强。这意味着脉冲不能无限早地布置下去——若种下得太早，它会在生长到有用规模之前就被黏性抹平；若种得太晚，则来不及长到需要的振幅。因此脉冲必须有一个精心安排的激活时刻，并由一个指数小的外力种子启动（手稿提到“指数小的外力种下每个脉冲”）。种子的量级必须与允许的生长时间相匹配：生长率与所需增益给定时间；时间给定允许的种子大小。

这解释了一个乍看奇怪的设计：脉冲要靠外力启动，而这个外力是构造本身安排的一环，且小到不影响外力整体的光滑性。在“残差博弈”的账本上，这是一笔划算的交易——用一个指数小的种子换来一个在奇点附近承担全部发散的正则化机制。

\subsection{离心不稳定性的物理判据}

离心不稳定性有明确的判据形式，可追溯到 Rayleigh 判据及其推广：在轴对称旋转流中，若角动量的径向分布向外递减得足够快（比 $r^{-2}$ 更快），则流是离心不稳定的。Leibovich--Stewartson 与 Billant--Gallaire 的准则把这套判据推广到非轴对称扰动。手稿安排核心外侧的角速度分布，使其恰好落在不稳定一侧；于是脉冲一旦被种下，其振幅增长由该判据保证的指数率驱动。

有一点值得注意：这里的稳定性判据与第 4 章的精确波振幅方程是同一件事的两种表述。Rayleigh 判据看的是角动量分布，振幅方程看的是 $D$ 的特征值——两者给出同一个“是否增长”的答案，只是语言不同。这种“同一事实、两种语言”的现象在本领域很常见，也是它难以初读的原因之一。

\subsection{尺度匹配}

脉冲的尺度必须与核心的坍缩尺度相配。核心以 $\ell_{r}\asymp\tau^{1/2}$、$\ell_{z}\asymp\tau^{1/2-h}$ 收缩，环域因此也在收缩；脉冲若要持续供给恰好抵消奇异残差的通量，其波长、振幅与激活时刻都必须随 $\tau$ 同步安排。手稿的作法是让脉冲“分族”布置——在固定高度上按整环排布，在径向—轴向上按局部小片排布——使得两类扰动分别承担径向与轴向的动量通量，且在给定的角向平均下相互补充。把这几套尺度对齐，是第 7 章逐阶修正能够启动的前提。

\subsection{与 Córdoba--Martínez-Zoroa 的差别}

同样是“被放大的扰动”，两处角色不同。Córdoba--Martínez-Zoroa 的放大扰动（涡层、振荡层）本身把解推向更奇异，奇点是放大的累积效果；OpenAI 的脉冲则服务于一个预先写定的光滑坍缩，其角色是把发散移交给 Reynolds 应力，从而让外力保持光滑。这个差别解释了为什么 \NS{} 的构造需要处理“脉冲—背景—外部流”三段式，而 IPM/Boussinesq 的构造更像是一串首尾相接的层。

\subsection{脉冲的几何：整环与局部小片}

手稿的脉冲不是各向同性的随机扰动，而是有明确几何。在固定高度上，脉冲包络铺满完整的环（其径向与角向扰动速度在环内变化），构成“整环族”；在径向—轴向截面上，脉冲则聚成一系列局部小片，每片是某个整环的横截面，另构成“小片族”。两族扰动的极化方向不同（一族给径向、一族给轴向），因此在给定的角向平均下能够互补：把它们的 Reynolds 应力 $\langle w\otimes w\rangle$ 加起来，正好覆盖背景残差需要被抵消的那几个分量。

这套几何安排有两个好处。其一是“平均为零、二次不零”这一性质得到保证：整环在小片尺度上仍然振荡，故其自作用的平均不为零，而它对背景的线性影响因平均为零而消失，从而只留下我们要的通量项。其二是与轴向外流的兼容：小片沿轴向梯度较大，可在轴向动量平衡中扮演主导，而整环在角向的变化则服务于角向平衡。附录 C 的“容许应力锥”就是把这套几何翻译成一个可以被逐点满足的应力实现问题——给定想实现的应力张量，问是否存在一族这样的脉冲去实现它，答案是肯定的（在一个锥内）。

\section{核心数学技术（三）：逐阶修正与外力光滑性}

\subsection{残差的阶次结构}

单个 ansatz 不可能让残差严格为零；它总会留下误差。这里的处置办法与第 3 章的纲领一致：把残差按阶次组织，用一个无穷序列的修正把误差不断推到更高阶，使其在每一个导数意义上都趋于零。手稿的第三至第十章就是这条修正链的展开，其骨架可以分成四段：构造主导流（第四章）、把基流修正到任意阶（第五章）、用辅助环面把振荡实现出来并修正残余应力（第六、七章）、以紧支撑的均值修正与残差改进收尾并推广到全空间（第八至十章）。

\subsection{与 Alpöge--Buckmaster 的机制对照}

Alpöge--Buckmaster 的 Boussinesq 构造里，有一个可以直接类比的机制：在包络变化之处，波 ansatz 会留下振荡误差，这些误差被显式修正到阶 $J_{q}=2k_{q}+8$，而 $k_{q}\to\infty$。让修正阶数随层号趋于无穷，是 Córdoba--Martínez-Zoroa 的“阶数随尺度上升”原则的延续，也正是外力级数能逐导数收敛的原因。\NS{} 的逐阶修正遵循同一逻辑，只是被修正的对象更复杂——除了背景残差，还包括脉冲之间、脉冲与背景之间的交叉项。

\subsection{构造链条与三则附录}

手稿把技术细节分摊到三个附录：附录 A 做径向矩匹配并构造热核型的“外部”，使核心能与一个速降的外部流对接；附录 B 给出轴附近的解析剖面及其延拓，为核心内层的正则性服务；附录 C 处理容许应力锥的实现，即把“想实现的 Reynolds 应力”翻译成“能实际布置出来的脉冲族”。这三块合起来，才使前面几章的标度与通量论证落地为一个真正全局定义的光滑构造。

\subsection{从主导流到全局构造：四步链条}

手稿第三至第十章的展开，可以归纳为一条四步链条，每一步解决一个具体障碍。

第一步（第四章）是先造出“在核心内精确成立、在外部光滑”的主导流：核心给出自相似坍缩，外部给出速降的缓流，两者在环域对接。这一步留下一个问题——环域的残差发散。

第二步（第五章）把基流修正到任意阶：以逐阶展开的方式，把残差里可消去的部分尽量在 ansatz 层面消去，使剩余误差的阶数尽量高。这一步的作用是降低后续修正的负担，但它无法处理环域中真正发散的那一块。

第三步（第六、七章）用辅助环面把振荡“实现”出来：先在一个模型环面 $\mathbb{T}^{2}$ 上把脉冲的构造与估计做干净，再把它搬回柱坐标下的环域，并用显式修正消除 ansatz 留下的振荡误差。引理 8.6 正是这一阶段反演方向导数算子的工具。

第四步（第八、九、十章）收尾：用紧支撑的均值修正把全局支撑压紧，改进残差的余项，最终把构造推广到整个空间并得到紧支外力。

这条链条的每一环都以“残余误差的阶数”为接口：前一步的输出是后一步的输入，接口上的阶数必须单调上升，这正是“修正阶数随层号趋于无穷”在本构造中的对应形态。读懂“阶数在链条上如何流动”，是读懂全篇的关键。

\subsection{外力为何光滑却不解析}

有一处细节必须在这里点明，因为它直接连到第 9 章的刚性定理。手稿的外力被定义为推进残差，它 $C^{\infty}$ 光滑，但并不是实解析的：一方面，它在空间上紧支撑，紧支撑的非零实解析函数不存在，故它在全空间上不可能解析；另一方面，它在奇点邻域也非解析，因为它由紧支撑的截断构造拼出。这个“光滑但不解析”的性质当时看起来只是技术副产品，事后却被证明是构造得以成立的必要条件之一——第 9 章会看到，若外力在奇点附近实解析，构造将被迫正则，与爆破相矛盾。

\section{核心数学技术（四）与验证争议：引理 8.6 的形式化落差}

\subsection{引理 8.6 的位置与作用}

引理 8.6 位于第八章“紧支撑均值修正”。它的作用是把一个局部问题降到环面 $\mathbb{T}^{2}=\mathbb{R}^{2}/\mathbb{Z}^{2}$ 上，并在那里反演一个方向导数算子。设 $F$ 是参数 $p$（外部空间变量）与环面变量 $y$ 的函数，固定方向 $v_{t}$ 定义算子 $N=v_{t}\cdot\nabla_{y}$。要在修正中“除以 $N$”解出某个势，就需要控制 $N^{-1}$ 对 $F$ 各阶导数的放大倍数。引理 8.6 给出的正是这样一个反演估计。它是整条修正链里一个不起眼却必须成立的环面估计。

\subsection{内联证明链：除子界、反演乘子与可和性}

在手稿的语境里，这个估计的证明很短，且每一步都可手推，值得完整走一遍，因为它把“数论里的除子界”和“二维 Fourier 级数的可和性”直接焊在了一起。

先看 Fourier 侧。把 $F$ 在环面 $\mathbb{T}^{2}$ 上展开，$\widehat{F}(k)$ 记频率 $k\in\mathbb{Z}^{2}$ 处的系数。算子 $N=v_{t}\cdot\nabla_{y}$ 在频率 $k$ 上作用为乘上 $i\,(v_{t}\cdot k)$，故其形式逆 $N^{-1}$ 的象征是 $\big(i\,(v_{t}\cdot k)\big)^{-1}$，放大倍数即
\begin{equation}
\big|v_{t}\cdot k\big|^{-1}.
\end{equation}
这是估计的第一步：只要方向 $v_{t}$ 固定且非零，频率 $k$ 离方向 $v_{t}$ 的垂直平面越近，$|v_{t}\cdot k|$ 就越小，放大就越厉害。手稿在此调用一个除子型的界，把这一放大用频率的总大小兜住：
\begin{equation}
\big|v_{t}\cdot k\big|^{-1}\;\le\;C\,(1+|k|).
\label{eq:divisor}
\end{equation}
式 \eqref{eq:divisor} 的含义是：最坏情形（$k$ 几乎垂直于 $v_{t}$）下，放大倍数被频率本身线性控制——因为整数格点靠近某个过原点的直线时，偏离量有其算术下界。

第二步是可和性。求导在频率侧对应乘上 $|k|$。设要求控制输出到 $m$ 阶导数，而输入比输出多付 $D$ 阶导数，则每个频率分量的贡献为
\begin{equation}
|k|^{m}\,\big|v_{t}\cdot k\big|^{-1}\,|k|^{-(m+D)}
\;\le\;C\,|k|^{m}(1+|k|)\,|k|^{-(m+D)}
\;\lesssim\;(1+|k|)^{1-D},
\label{eq:summ}
\end{equation}
其中第二步用了除子界 \eqref{eq:divisor}。于是问题归结为：$D$ 取多少，才能让 $(1+|k|)^{1-D}$ 在二维格点上可和。二维的判据是
\begin{equation}
\sum_{k\in\mathbb{Z}^{2}}(1+|k|)^{-s}<+\infty\iff s>2,
\end{equation}
因为半径 $|k|\sim R$ 的格点数按 $R$ 增长，与 $(1+|k|)^{-s}$ 相乘得到 $\sim R^{1-s}$，而 $\sum_{R}R^{1-s}$ 收敛当且仅当 $1-s<-1$，即 $s>2$。记衰减式 $(1+|k|)^{1-D}=(1+|k|)^{-s}$ 中的 $s=D-1$，则二维可和要求 $s>2$，即 $D>3$；整数下最小可行值即为 $D=4$，此时衰减恰为 $(1+|k|)^{-3}$，可和，但距不可和的 $(1+|k|)^{-2}$ 只差一阶、没有任何余量。这就是“比输出多四个输入导数”的来源。若换到三维环面，判据变为 $s>3$，即 $D>4$，最小整数为 $D=5$——所以这条引理必须在二维环面上陈述，多一维就不够。

至此可和性成立，进而给出估计 \eqref{eq:summ} 对应的算子界。手稿随后补一句：该乘子与一切“慢变量”系数导数可交换，且保持函数的实性；再补一句：固定零频系数后解唯一。这三句话合起来，就是引理 8.6 的完整陈述。

\begin{remark}
这套证明有一个容易被忽略的结构特征：唯一真正用到“数论”的地方是除子界 \eqref{eq:divisor}，其余部分只是 Fourier 级数的初等记账。之所以要绕道频率侧，是因为方向导数算子 $N=v_{t}\cdot\nabla_{y}$ 在物理侧不可逆（沿 $v_{t}$ 方向为常数、沿垂直方向属零空间），而在频率侧它的逆只是一个逐频相乘的因子；代价是这个因子在 $v_{t}\cdot k$ 接近零时爆炸，需要除子界来控制。这类“在频率侧求逆、用算术界兜底”的手法，属于拟微分算子记帐的常见套路，但它在这里被安排在一条纯初等的证明里，用起来很轻。

还须补上关于这条界本身的一点：它不是自动成立的记账琐事，而是一条带算术内容的输入。对有理方向，$|v_{t}\cdot k|^{-1}$ 一致有界，远强于所需；对无理方向，Dirichlet 逼近定理保证 $|v_{t}\cdot k|\lesssim|k|^{-1}$ 对无穷多个 $k$ 成立，最坏情形下增益因子 $(1+|k|)$ 被用满——这正是 $D=3$ 不足、$D\ge4$ 的原因；若方向具有更极端的有理逼近（Liouville 型），这条界本身会失效，任何固定导数额度都不足以收尾。因此 $m+4$ 不能被视为记账的自动推论：它依赖方向 $v_{t}$ 的算术性质，而该性质在构造里由对方向的选择决定。
\end{remark}

\begin{keybox}
\textbf{证明链的三步}\quad （1）逆算子的象征是 $|v_{t}\cdot k|^{-1}$；（2）除子界把它换成 $C(1+|k|)$；（3）输入比输出多给 $D$ 阶导数后，衰减为 $(1+|k|)^{1-D}$。二维格点的可和判据是 $s>2$，故 $D=4$ 时给 $(1+|k|)^{-3}$（$s=3$，恰可和、无余量），这正是“多四阶”的来源；三维判据是 $s>3$，需 $D\ge5$，所以引理只能在二维环面上陈述。
\end{keybox}

\subsection{落差：多四个导数还是多五个}

争议出现在自然语言证明与 Lean 形式化之间。手稿正文在引理 8.6 的开头给出如下论证：由除子界 $|v_{t}\cdot k|^{-1}\le C(1+|k|)$，对输入 $F$ 求的导数比所求输出多四个，即留下二维可和的 $(1+|k|)^{-3}$ 界，式 (8.19) 由此得证。但 Bastounis、Circelli 与 Hansen\cite{bch}指出，Lean 仓库中与式 (8.19) 最接近的定理 \texttt{norm\_derivativeWord\_inverse\_le}，其假设用的是“多五个导数”，即控制同一批环面输出导数时，Lean 版本要求输入多出 $m+5$ 阶而非 $m+4$ 阶。就命题强度而言，多要求一阶输入正则性是更强的假设，因此 Lean 该定理所证明的命题弱于手稿声称的命题；这正是研究者所称“自然语言论文声明了一个更强的反演估计”的所指。

联系上一小节的记账式 \eqref{eq:summ} 可以把这一差别算准：该式对一般的输入导数额度 $D$ 给出衰减 $(1+|k|)^{1-D}$，与 $D$ 取四或五无关地成立。代入 $D=4$ 得 $(1+|k|)^{-3}$（$s=3>2$，可和但无余量），代入 $D=5$ 得 $(1+|k|)^{-4}$（$s=4>2$，可和且有余量）；两者在二维环面上都收敛；不可和的只有 $D\le3$，其中 $D=3$ 恰为临界情形（衰减 $(1+|k|)^{-2}$）。因此落差之处不在可和性，而在命题的强弱：$m+4$ 是更强的估计，它用较少的输入正则性控制同一批输出导数；$m+5$ 则较弱。手稿第 96 页声称的 (8.19) 属于前者，而 Lean 中对应的定理证明的是后者。按 Bastounis 等人的表述，形式化在这里并未覆盖手稿声称的那条更强的估计。

以上落差此前只见于 Bastounis 等人的转述。本文进一步核对了公开仓库的一手源码，结论与转述一致。其一，仓库 \texttt{openai/}\allowbreak\texttt{NavierStokesAndEuler} 在 commit \texttt{f9e8bc5} 处的完整文件树共 2673 个条目，其中 \texttt{.lean} 文件 2659 个；在该仓库内检索声明名 \texttt{norm\_derivativeWord\_}\allowbreak\texttt{inverse\_le}，命中文件唯一，即 \texttt{NavierStokes/}\allowbreak\texttt{SmoothFamilyTorusInverse.lean}。其二，该文件共 1140 行、57334 字节，blob 哈希为 \texttt{c607eb6f2460}\allowbreak\texttt{a3ba55dc2218}\allowbreak\texttt{072cd504bf62}\allowbreak\texttt{53bf}，被质疑的定理声明位于第 1059 行：其两个假设项分别对 $F$ 的切片与交换切片要求 \texttt{xJet (w.length + 5)} 有界，即输入比输出多五个导数；结论侧的权重因子写作 $(\texttt{weight}\ k^{4})^{-1}$，指数为 4。其三，向下游看，同一文件第 1082--1136 行的 \texttt{mixedJet\_}\allowbreak\texttt{inverse\_bound} 在该段内共十处沿用 \texttt{w.length + 5}，并在第 1135 行调用上述声明；可见 $m+5$ 的口径并非孤立一处，而是贯穿了这条估计链。核验所得的代码与 Bastounis 等人所引代码块逐字一致（对方仅另加了红色标注）。因此落差的 Lean 一侧——形式化定理确按 $m+5$ 表述——是经一手源码确认的事实，而非仅凭二手转述；正文一侧的 $m+4$ 记述则依据 Bastounis 等人对手稿原文的转引。

由此产生的是一个具体的问题清单：其一，(8.19) 在其字面强度（$m+4$）下是否成立；其二，若它只以 $m+5$ 成立，那么修正链中按 $m+4$ 的强度调用它的位置，其输入函数是否恰好备有多余一阶的正则性可供支取；其三，Lean 定理的结论侧是否与 (8.19) 逐字对应——Bastounis 等人的表述是“控制同一批环面输出导数”，但输出界中常数的具体形式是否与正文一致，则仍需逐式对照原文确认（代码一侧的假设强度与权重指数已由本文核验，见上段）。需要说明的是，这处落差表明的是“形式化覆盖的命题弱于正文声称的命题”，并不等同于正文的 (8.19) 有误：按第 8.2 节复原的四阶记账，$m+4$ 的论证在简化模型中是自洽的，落差的实际后果取决于调用点处是否留有余量。三问都需要逐处核对正文与该引理的调用点，本文无法依据公开材料给出结论。

\subsection{SCI 谱系论证：为何“翻译”可能比停机问题更难}

Bastounis--Circelli--Hansen 的更大主张不止于这一个引理。他们区分两种自动形式化：第一种是把一份“自然语言证明”翻译成 Lean 并让 Lean 通过编译；第二种是把一份数学文本“语义忠实地”翻译成 Lean。他们指出，第一种的成功不能推出原自然语言证明正确，因为 AI 可以在翻译时悄然改换语义——或者给错误证明配上一个正确证明，或者把正确证明换成另一个正确证明，或者像本案这样，把“更强的声明”翻译成“更弱的定理”而仍然编过。

他们把语义忠实翻译的困难度放在可解性复杂度层级（Solvability Complexity Index，SCI）上：要消除自然语言数学文本的歧义（这是语义忠实翻译的必要条件），其难度在 SCI 谱系中高达无穷（$\mathrm{SCI}=\infty$），用他们的话说，比包括停机问题（$\mathrm{SCI}=1$）在内的任何计算问题都更难。由此得出的操作性结论是：OpenAI 的自然语言证明与其他被自动形式化的证明，不应仅凭 Lean 证书就获得信任，而应经受与其他证明同样的同行评审与审视。

评估这道落差时需要同时握住两侧的事实。按 OpenAI 公开的 Lean 仓库及第三方对仓库的独立计数，形式化包含六十余万行 Lean 代码、三万余条定理与引理，未留未证明的 \texttt{sorry} 占位，也未引入标准三公理之外的公理；但仓库自述的评审状态是“self-assessed（自评）”。Bastounis、Circelli 与 Hansen 强调，自评状态与独立验证是两回事，而引理 8.6 的样本说明：即便 Lean 全部编译通过，正文论证的某些步骤仍可能未被覆盖\cite{bch}。两方援引的都是公开可查的材料，分歧的解决有赖于逐条核对正文与形式化的对应关系；本文对此不作裁断。

\subsection{这一落差是否影响主定理}

诚实的回答是：目前无法从公开材料确定。需要区分两个层次的问题。其一是“引理 8.6 按其字面（$m+4$）是否成立”；其二是“即便该引理应按 $m+5$ 陈述，主定理的证明链是否会因此断裂”。第二种情形的可能性是存在的：一个中间引理的常数或导数阶数，未必是主定理链条上不可替换的一环——若其他步骤留有冗余，或该引理可以换用别的估计，主定理仍可能成立。反之，若该引理恰是拼接处唯一的支点，落差就会传导到主定理。这个判断需要把引理 8.6 在整条修正链中的调用点逐一梳理，超出本文范围，留作开放问题（见第 11 章）。

从方法论上说，这件事给出的教训比它给出的结论更重要。自然语言证明的检验，历来靠“读者复算”；形式化证明的检验，靠内核编译。当两者之间存在一层 AI 自动翻译时，“复算”与“编译”都不再直接作用于原文——这就是 Bastounis 等人所说的断层。对本领域而言，一个可行的折中是把形式化的对象从“全文”缩小到“关键节点”：对每个支撑主定理的核心引理，独立写一份与正文严格对应的形式化陈述并加验收，而不是依赖一次性把整篇论文译成代码。第 11 章的开放问题会把这一条列为具体建议。

% ============ 第三部分：第 9--12 章、附录与参考文献 ============

\section{数学界的验证与约束}

前三章处理的是“构造是怎么做的”。这一章处理“别人怎么看”，并特别关注一条容易被忽略的区别：验证文献里既有\emph{否证性}的结论，也有\emph{约束性}的结论，二者对构造的意义完全不同。

\subsection{Tao 的独立评估与方法溯源}

Terence Tao 在 9 月 7 日的博文\cite{tao-blog-sep7}中给出了一份独立于两家机构的评估。他称 Alpöge--Buckmaster 的结果是“了不起的成就”，把基本路线的归属明确记在 Córdoba 与 Martínez-Zoroa 名下，并指出：这条路从“外力下爆破”延伸到 \NS{} 是原则上可预期的，甚至有机会去掉外力。值得注意的是，这篇博文写于 OpenAI 公布 \NS{} 结果之前——也就意味着，“从 Alpöge--Buckmaster 的强迫 Euler 到千禧年问题，道路对任何理解这套方法的人都是可见的”这一判断，是一份独立的、时序在先的材料。

Tao 同时强调，他仍需要更多时间来消化细节，并把“突出这些工作的核心想法”视为这类评述的主要价值——求解问题只是发展数学理解的代理目标，而非目的本身。这句话为第 10 章的认识论讨论定了调。

\subsection{Constantin--Ignatova--Vicol 的刚性定理}

9 月 17 日，Peter Constantin、Mihaela Ignatova 与 Vlad Vicol 上传了一篇预印本\cite{civ}，其结论初看像是冲着 OpenAI 的构造去的，实则需要仔细分辨方向。

他们的定理说：设 $(u,\pi)$ 是带外力的三维 \NS{} 在单位抛物柱上的适当弱解，在紧子集上光滑；设外力在奇点时刻前一致地有 $C^{2}$ 界，且在空间变量上局部一致实解析；再设速度满足两条与手稿构造相对应的性质——(i) 其角向平均 $\mathcal{P}u$ 满足各向异性的 Type II 界（$v_{r}$ 被 $\tau^{-1/2}$ 控制，$v_{z},v_{\theta}$ 被 $\tau^{-1/2-h}$ 控制，每求一次 $r$ 导数付 $\tau^{-1/2}$、每求一次 $z$ 导数付 $\tau^{-1/2+h}$），(ii) 在一个坍缩的核心球上速度严格轴对称——那么该假想奇点实际上是正则点。

换言之：在这两条性质加上解析外力成立时，奇点不可能出现。由于 OpenAI 的构造恰恰满足 (i) 与 (ii)，定理的推论就变成一句约束：任何具有性质 (i)、(ii) 且外力在奇点时刻前保持 $C^{2}$ 有界的构造，其外力既不能在奇点附近恒为零，也不能在空间变量上实解析。作者还特意在手稿的附录里核对了 OpenAI 构造确实满足 (i) 与 (ii)，并把对方的外力性质定为“$C^{\infty}$ 光滑但非实解析”。

\subsection{关键：这是约束，不是否证}

初学者容易把这条定理误读为“OpenAI 的构造被证伪了”。恰恰相反。

手稿的外力由紧支撑的截断拼出，因而是光滑但\textbf{非}解析的（第 7 章已说明：紧支撑的非零实解析函数不存在，奇点邻域也非解析）。刚性定理说的是“若外力解析则正则”，而构造的外力不解析，因此它落在定理假设之外——定理对它不作任何断言，更谈不上否定。定理真正的价值在三处：其一，它说明了构造中哪些特征是\textbf{本质}的——各向异性 Type II 标度、坍缩核心的精确轴对称、以及外力不解析，这三者缺一不可，否则奇点会被刚性逼回正则；其二，它给出了一个可判定的“必要形态清单”，为后来者提供了检查表；其三，作者明确声明“我们并不声称验证了手稿构造的正确性”，只是从手稿陈述中抽取所需性质。这种克制本身就值得记录。

\begin{principlebox}
\textbf{约束性结论的正确读法}\quad “在更强假设下结论为正则”，并不等于“原构造错误”，而是给出“原构造必须避开这些假设”这一必要条件。CIV 的贡献在于把“外力必须非解析”从技术细节升格为结构必然。
\end{principlebox}

\subsection{刚性定理的证明思路：放大尺度与古代解}

CIV 的结论不只是一条约束，它的论证方式本身值得一记，因为它示范了当代正则性理论的一种标准操作。思路分三步。

第一步是放大。利用各向异性的 Type II 界所提供的长度尺度 $\ell_{r}(t)=(-t)^{1/2}$、$\ell_{z}(t)=(-t)^{1/2-h}$，在假想奇点处做抛物型伸缩。这一步的关键在于：Type II 标度允许把速度场规范化，使伸缩后的解在紧集上一致有界——否则连取极限都无从谈起。

第二步是取“古代解”。伸缩参数趋于零时，得到的极限解定义在 $(-\infty,0]$ 上，即一个自古以来就存在的解。这类对象的约束极强：它没有初始时刻，因而任何依赖“信息从初值单向流入”的论证都被排除。

第三步是刚性。对古代解施加条件（外力的解析性、解的轴对称性），证明这样的古代解必然平凡，从而把假想奇点排除。作者指出，这一步的机制与他们此前关于轴对称 \NS{} 的工作一脉相承。

这三步合起来说明：刚性定理真正用到的是假想奇点附近解的形态足够特殊这一事实；它并未用到被质疑的引理 8.6。两条验证线互不干涉，这也是本文把它们分开处理的理由。

\subsection{Clay 的程序}

克莱数学研究所 9 月 11 日的表态\cite{clay-news}只说问题“看上去已被解决”，并强调评估与授信的流程“刻意不匆忙”。按其规则，授奖需要满足若干条件，包括在合格期刊发表、发表后至少两年、以及数学界的普遍接受；公开材料显示至今没有授奖决定，且 OpenAI 表示不打算为此申领那笔奖金。把这三件事放在一起看，可以得到一个平稳的结论：Clay 的表态是“注意到了一项重要声明”，不是“认证了一个定理”。

\subsection{形式化验证能保证什么、不能保证什么}

Lean 证书给出的是：某个命题在一套公理系统内可证，且证明项合法。它不能自动给出的是：这个形式化命题与手稿中那句自然语言命题是否\textbf{同义}。二者之间的空隙，正是第 8 章所述争议的所在。对此，学界有一句被反复引用的比方——耶胡达·拉夫（Yehuda Rav）说“定理是标题，证明是正文”：一个正确但未被消化的证明，只给了标题而没给正文。这句话出自 Tapio Schneider 发表于 Tao 博客的客座文章\cite{schneider}，它把讨论从“对错之争”引向“理解之争”，也正是下一章的主题。

\section{认识论与学术生态：形式化、理解与归责}

\subsection{形式化不等于理解}

把前几章的线索接起来，可以得到一个不依赖任何争议事实的判断：形式化验证与数学理解是两件不同的事。形式化解决的是“从公理到命题的推理是否合法”，理解解决的是“人为什么能看出这个命题应当成立、这套方法能迁移到哪里”。前者可以自动化，后者目前不能。Schneider 的文章提议把 AI 放进“可审计的脚手架”里使用——在数学里是 Lean，在物理与工程里是高分辨率模拟、观测或实验——让 AI 在每一环都可被单独检验的位置上发挥作用。他的例子很干净：天气预报这类每天都能核对的预测，黑箱模型可以独立成立；而气候外推、飞行器定型这类预测必须在被经验检验之前就获得信任，此时可审计的因果链才是信任的来源。

值得补一句的是：形式化与理解的分野，在数学内部早已有先例，并非 AI 带来的新问题。一个极端的例子是四色定理：1976 年由计算机辅助完成，其论证此后被部分学者质疑过“是否算证明”，直到后来的形式化工作把可核验的部分做得更彻底。另一个例子是有限单群分类：其证明长达上万页，几乎没有人从头到尾读完，社区接受它靠的是分层信任——大量中间结果被反复重证与引用。数学共同体早就活在“不完全理解却接受结论”的状态里；AI 只是把这种状态从“少数巨型证明”扩展到了更多场合，并让它变得廉价而频繁。

\subsection{理解的代价：为什么“消化”如此重要}

还有一层更温和但同样重要的理由，说明为什么“可理解”不只是审美要求。数学知识的再生，靠的是方法在人群之间的传播：一个被理解的构造会衍生出新的想法、新的练习、新的应用；一个只被机器检验过的构造则停在那里，除非有人把它想通。Rav 的“标题与正文”是为了说明这一点，Schneider 的 episteme/techne 之辨也是为了说明这一点。

具体到本项工作，“消化”的难度是肉眼可见的：166 页正文加两百余万行 Lean；外部读者若要完整核对，需要同时具备调和分析、流体方程与形式化工具三方面的能力。这不是 OpenAI 独有的问题——Buckmaster 也坦承自己那篇 Euler 写稿“像 AI 生成垃圾”，并为此道歉。当构造的复杂度超过共同体的消化带宽时，问题不是“对不对”，而是“多久以后才能真正被人用起来”。

\begin{quotebox}
“定理是标题，证明是正文。”——一个正确却不被理解的证明，给出了标题，却没有交代清楚证据链。Schneider 的论点，是把 AI 安排在证据链上那些能够被单独检验的环节，而不是让它端到端地给出结论。
\end{quotebox}

\subsection{episteme 与 techne}

Tao 与 Schneider 都借用了古希腊的一对概念：episteme（解释性理解）与 techne（预测与制作的技艺）。求解一个方程的解的性态，属于 episteme；让模型预报明天的天气，属于 techne。Schneider 的观察是，当预测能被廉价地反复核验时，techne 可以独立成立（黑箱天气预报就是例子）；当预测必须在被核验之前就被信任时，episteme 与 techne 便不可分离，可审计的证据链随之成为硬约束。这一段论证对数学的启示是：Lean 让它变成“可以被机器核验的 techne”，但数学共同体真正消费的仍是 episteme——这解释了为什么一个 Lean 编过的证明，仍可能被数学界抱怨“没有交付理解”。

\subsection{优先权、训练数据与署名的争议记录}

围绕本项工作产生的一组事实主张，涉及研究生态与科研伦理，而不涉及数学内容本身。以下叙述仅依据公开材料并逐一标明来源。这些主张之间存在直接冲突，本文不掌握足以裁断的材料，因而不对任何一方的主张作出裁断。

关于优先权，据 Buckmaster 的公开声明\cite{buckmaster-statement}，他与 Alpöge 在 8 月中旬已获得光滑外力下 IPM、Boussinesq 与三维 Euler 方程的爆破结果，并认为自己同时握有次耗散 \NS{} 的结果，只是由于 Lean 验证尚未完成而未一并发布。OpenAI 的官方说明\cite{openai-blog}则称，其相关工作自 9 月 1 日因听闻“两个千禧年问题已被解决”的传闻而启动；公司表示其研究人员与智能体在对方公开发布之前没有看到对方的任何工作，并说明在 9 月 6 日完成 Lean 验证后曾联系对方，提出同步发布、承认对方在外力 Euler 上的优先权。双方对若干关键时间节点的叙述存在出入，特别是首个提示词的发出时间；在缺乏版本历史等客观记录的情况下，这一争议无法由外部材料裁决。

关于训练数据，据 Buckmaster 的声明\cite{buckmaster-statement}，他曾询问其与 Alpöge 存放在 Codex 会话中的草稿是否被用于训练或检索，未获得直接答复。OpenAI 方面表示，其模型没有查询特定用户数据，同时承认“虽不太可能，但不能排除来自用户使用的去标识化数据对模型改进有所贡献”。围绕这一表述，后续的技术讨论提出了一种更窄、更可检验的解读：相关机制更可能是对未选择退出的用户对话进行强化学习，而非直接“查阅草稿”；真正待定的问题因而变为最近一次训练发生的时间，以及模型是否会记住强度足以影响结论的片段。这一表述与 OpenAI 的措辞相容——公司否认访问特定用户数据，但不排除去标识化的派生数据对模型有所贡献。

关于署名，手稿封面唯一的署名是“OPENAI”，这一机构署名不区分研究纲领的来源、中间推进、\NS{} 特有的构造与形式化工作。对此的批评集中在责任与归责上：当外部研究者质疑某个引理或形式化对应关系时，由谁出面回应并不明确。支持机构署名的常见理由则是：成果确系组织内部的系统与团队共同产生，个体贡献难以按传统学术出版的方式切分。这一分歧是学术出版规范在面对新型研究组织形态时的一个具体案例。

关于通话内容的争议，据 Buckmaster 公布的叙述，他在 9 月 6 日与 OpenAI 方的通话中被提出两种发表方案，并记录了若干他本人视为施压的话语。被点名的 OpenAI 方面研究者公开称这些指控“虚假且具有煽动性”，表示自己是遵循学术规范参与讨论，并承诺给出更完整的说明。同一事件的两种叙述就此直接冲突。

需要区分的是“数学有效性”与“产生条件是否正当”这两件事：一个证明可以在产生条件存疑的情况下依然正确；反过来，产生条件的争议也不构成对该定理本身的数学反驳。第 9 章的约束性结果属于数学层面；本节的记录属于事实主张层面，且这些主张的来源、口径与可核验程度各不相同，读者宜按上文的归因分别看待。

\subsection{七百余篇手稿与负责任披露之争}

10 月的报道称，OpenAI 已放出数百篇 AI 生成的数学手稿，并与学界就披露方式产生摩擦：一些参与过 8 月咨询的数学家反映，他们希望以正式论文而非博客或社交媒体形式披露的意见未被采纳。把这些放进本文的框架里看，它是第 10 章那组张力的一次放大——当产出速度远超共同体消化速度时，验证、归属与理解三件事会同时承压。本节仅记录事实，不作价值判断。

\subsection{小结}

把这一章压成一句：\NS{} 事件的数学部分是清晰的（一条被验证过方向、被约束过形态的构造路线），它的社会学部分则远未尘埃落定。对本文的读者而言，重要的能力是能同时握住这两条线，而不把它们混为一谈。

\section{诚实的边界与开放问题}

\subsection{备选 A、B 仍未决}

所有公开材料一致指出：本次结果是备选 C 与 D，无外力的 A 与 B 依然开放。9 月 20 日，Petrillo 与 Glimm\cite{petrillo-glimm}提出了“正缺陷问题”作为无外力搜索的靶标：要求周期盒上从光滑数据出发的 Leray--Hopf 解在某个有限时间窗内损失的能量，超出黏性所耗散的部分。正缺陷蕴含爆破，因而蕴含备选 B 的否定；其逆向是否成立尚不清楚。他们把该靶标等价地化为 Littlewood--Paley 壳层能量通量的下界（即 Onsager 湍流理论中粗粒化通量的 Fourier 侧形式在窗口上的平均），并列出候选解必须满足的必要条件。这项工作本身也在 Lean 4 中形式化，且刻意在假设里引入 \NS{} 对象，使方程只通过假设进入。它对本文的意义是：它把“无外力路径到哪一步、缺什么”讲清楚了——缺的不是方法，而是一个靶标。

\subsection{解析外力障碍：可绕开还是本质}

第 9 章的刚性定理把外力非解析变成了必要条件。接下来的问题是：这一障碍是构造的偶然特征，还是 \NS{} 强迫爆破的固有属性。若它是固有属性，那么“光滑外力”这条路本身就带着一条暗线约束——外力必须在奇点附近保持非解析；若它只是当前构造的副产品，未来或许能构造出外力更“规整”的爆破。公开材料目前不足以回答，但它给出了一个明确的检验方向：尝试在保持 (i)(ii) 的前提下改进外力的解析性，看刚性定理在哪一点被触发。

\subsection{引理 8.6 落差的影响范围}

如第 8.5 节所述，这一落差是否传导到主定理，取决于该引理在修正链中的调用结构与余量。落差的“Lean 一侧”已无争议：本文据公开仓库 commit \texttt{f9e8bc5} 一手读取了该声明（文件 \texttt{NavierStokes/SmoothFamilyTorusInverse.lean} 第 1059 行，两个假设项取 \texttt{w.length + 5}，权重指数为 4，详见第 8.3 节），确认其字面强度确为 $m+5$；于是争点不在“Lean 写了什么”，而在“$m+4$ 的正文声明是否成立、以及在何处被调用”。一个可操作的收敛问题是：把引理 8.6 从 $m+4$ 改为 $m+5$ 后，其后各章中依赖它的估计是否仍成立；若成立，则落差是局部可修的；若不成立，则需在更上游重做。这需要逐条比对正文与 Lean 的对应关系，属于可以完成但颇为费时的核查，也是验证工作最应当优先投入的地方。

\subsection{构造能否被简化到人类可消化}

166 页正文加两百余万行 Lean，其规模本身就是理解的门槛。历史上，重大构造通常要经过数年乃至数十年的简化、重讲，才会进入教科书。是否已有或将会出现一份把核心思想压缩到数十页、可被研究生读懂的阐述，是判断这项结果能否真正进入数学共同体的关键。已有的“友好读本”（如 Shavit 的机制笔记\cite{shavit}）是这类努力的雏形，但它们解释机制，尚未重构证明。

\subsection{反方证据：把“解决”读成“终局”会错在哪里}

一篇诚实的技术综述必须把反方证据摊开。反对把此事称为“千禧年问题已解决”的理由至少有六条，它们彼此独立，且都建立在公开材料之上。

第一，方向性。结果落入备选 C、D（有外力），而最常被公众理解为“水的方程”的备选 A、B（无外力）仍然开放。用“有外力则可爆破”去回答“无外力是否永远光滑”，是换了一个问题。

第二，验证层次。Lean 通过只保证“形式命题可证”，不保证“形式命题与原命题同义”。Bastounis 等人已在引理 8.6 处给出一个具体的不一致样本。真正的独立复核（由第三方重做关键估计）尚未完成。

第三，传出规模。手稿 166 页、Lean 两百余万行，共同体的消化需要时间；Clay 的规则也明确要求“两年”与“普遍接受”。9 月的表态不是认证。

第四，构造的人为性。这是一条被精心设计的路径，依赖大量参数选择；它不给出“一般流会怎样”的刻画，因而对理解真实湍流（Kolmogorov 尺度以上的、普遍存在的湍流）帮助有限——Schneider 的文章专门指出了这一点：该结果“对理解物理湍流几乎没有发言权”。

第五，归属与条件争议。优先级、训练数据、署名方式目前均未定论，且两方叙述直接冲突。这不是数学错误，但足以说明“共同体的接受”还没有发生。

第六，机制而非结论的进展。CIV 的刚性定理表明，任何满足同样形态的构造都必须使用非解析外力。这既是重要洞察，也提示：这条路径可能带有一条内在的形态约束，而不是一条通用的爆破机制。

\begin{keybox}
\textbf{平衡读法}\quad 把这些反方理由与第 13 章的四点定位并置，得到的不是“模棱两可”，而是一个精确坐标：数学结论在特定范围内成立且已给出形式化证书（自评）；无外力情形未决；独立验证与共同体接受尚在途；物理意义有限；归属未定。把任何一项单独放大，都会得到失真的结论。
\end{keybox}

\subsection{可检验的收敛性建议：尚未完成的严格推演}

在给出方向建议之前，需要把“尚未完成”这件事说准。本项工作并非没有形式化——OpenAI 公开了包含六十余万行 Lean、三万余条定理与引理的仓库，且未留未证明的占位、未引入标准三公理之外的公理。真正尚未完成的是三件事：其一，该证书的\textbf{独立核验}（仓库自述评审状态为“自评”，未见第三方逐条复核的公开结果）；其二，证书与正文之间的\textbf{语义忠实性}（引理 8.6 处已有具体反例）；其三，其\textbf{数学思想的严格重述与推演}——即由共同体独立地、以可读的方式把“为什么这套构造能成立”重新推导一遍。第三件事之所以关键，是因为在上述两项完成之前，我们尚不能确定自己讨论的究竟是一个已确立的定理，还是一个有待确立的候选定理；而第 12 章关于“思想定位”的判断，恰恰以“我们确实理解了证明的动作”为前提。

以下任务按“先证成、再消化、后拓边”的顺序排布，供后来者按序推进。

其一，\textbf{核对引理 8.6}。对照手稿正文第 96 页附近的 (8.19) 与 Lean 仓库 commit \texttt{f9e8bc5} 中 \texttt{norm\_derivativeWord\_inverse\_le} 的假设，判定字面强度（$m+4$）是否为真；若为真，补一份与正文严格对应的形式化陈述；若为假，向上游追踪该引理被调用的位置，判断主定理是否仍成立。

其二，\textbf{复核刚性定理与技术细节的对应}。逐条核对 CIV 附录 A 中从手稿陈述抽取的性质 (i)(ii)，确认它们确实是手稿所证，而非对手稿的额外假设。

其三，\textbf{语义忠实形式化的协议}。对每一个承重引理，独立撰写一份与正文一一对应的形式化陈述，并附“对应性证据”（逐行或逐式对照说明），而不是依赖一次性把整篇论文译成代码。这是把第 8、9 章所述断层缩小的最直接手段。

其四，\textbf{构造的可简化}。尝试把“脉冲供给 Reynolds 应力”这一机制从 166 页的估计中剥离出来，写一份数十页、可被研究生独立读通的重述。

其五，\textbf{加强结论}。审视备选 C 是否可加强到“外力实解析”或“外力可取零”的变体；这一步直接与 CIV 的障碍对话，也是通往备选 A、B 的必经关口。

其六，\textbf{机制的可迁移性}。把“零平均振荡经二次自作用生成指定平均通量”抽成一个可复用的模块，测试其能否迁移到磁流体方程、次耗散 \NS{} 或其他流体模型；能否迁移，是判断它属“纲领内一步棋”还是“可推广的新工具”的实测标准。

其七，\textbf{结构分类}。以 CIV 型刚性定理为工具，刻画“哪些（标度、对称性、解析性）组合允许爆破、哪些被刚性排除”，把当前构造放进一张完整的分类表中；这既是验证工作，也可能产出独立的新定理。

其八，\textbf{无外力靶标的检验}。检验“正缺陷问题”作为无外力靶标的可行性：Petrillo--Glimm 已把靶标化为 Littlewood--Paley 壳层通量的下界，并做了 $128^{3}$、$256^{3}$ 的伪谱搜索；需要进一步判断该路径能否导向证明。

\subsection{这件事与“AI 能不能做数学”是两问}

最容易发生的误读，是把本项工作直接读成对“AI 是否会做数学”的回答。这两件事需要分开。

严格地说，本项工作证明的是一个数学命题（在其成立的前提下），而不是一个关于 AI 能力的定理。三点理由。其一，过程有大量人类参与：OpenAI 的团队设计了多智能体框架、划分了问题变体、在过程中更新模型、用 Codex 汇总各组洞见；按其自述，问题变体（A/B/C/D 分发给不同组）这一步是人安排的。其二，路线有明确的人为先驱：Córdoba--Martínez-Zoroa 的纲领、Alpöge--Buckmaster 的推进，都是人类在 AI 之前数月乃至数年铺设的。其三，验证环节依赖的是 Lean 内核，一种与人无关但对“理解”也无关的形式工具。

反过来说，也不能因为“有人类参与”就否认这次跃迁。把一份“历史上最糟糕的写稿”在七天内变成 Lean 可验证的证明，这件事本身在几年前是不可想象的；它确实标记了工具能力的台阶。诚实的立场是：这是一次“人设定方向与判据、机器承担搜索与记账”的协作突破，其意义在于改变了求解环节的成本结构，而不在于宣告某个关于智能的命题。

\subsection{长期含义}

把视野拉长，这件事提出的问题比它回答的问题更多。当“求解”的速度远超“消化”的速度时，数学共同体的训练方式、评审方式、归功方式都面临重估。Tao 的立场是：理解与求解是光谱的两端，有人以求知为目的，有人以求解为目的，二者都值得保留，不应互相规定何者为“正”。Gowers\cite{gowers}则更直接地把这称为一次“苦涩的教训”——他所在的自动定理证明小组试图通过理解人类如何证题来设计算法，而大模型绕开了这种理解照样做到了。这一判断与本文的技术观察一致：工具层面的自动化已经发生，而解释层面的自动化仍未到来。

\section{数学思想的定位：开创性突破，还是既有纲领的凸闭包？}

\subsection{问题的提法}

一个自然的问题是：OpenAI 在 \NS{} 上的工作，究竟发明了新的数学思想与方法，还是既有数学知识与方法的闭包推论？换言之，它是否带来了数学思想层面的革命性创新？

要回答，先得把“一个结果在数学思想史中的位置”拆成可以分别判断的层次，否则“创新”一词会在不同尺度上滑动。本章采用三层判据：\textbf{纲领层}，指研究纲领、问题的重新表述与核心图景；\textbf{技术层}，指为跨越具体障碍而引入的构造步骤与估计；\textbf{方法论层}，指生产与验证知识的流程与工具。一个工作可以在某一层是革命性的，而在另一层只是收官；把三层混为一谈，是这类讨论中最常见的失误。

\subsection{纲领层：不是它发明的}

判据一：核心图景是否出自本工作。答案是否。

“先选解、后定力”的残差博弈、逐尺度放大、修正阶数随尺度上升，这三件事构成整个工作的思想内核，而它们的归属在公开材料中高度一致。Buckmaster 在声明中把纲领的基本想法记在 Córdoba 与 Martínez-Zoroa 名下，并称该攻击路线“不是由我们提出，也不是由大语言模型提出”；Tao 在 9 月 7 日（早于 OpenAI 公布结果）的博文中独立地给出同一条溯源，并明确指出从 Alpöge--Buckmaster 的强迫 Euler 到 \NS{} 的道路“对任何理解这套方法的人都是可见的”。按此，纲领层不属于 OpenAI 的原创。

需要谨慎的一点是：可预期不等于可推出。一个结果落在既有纲领的射程内，并不自动意味着它是该纲领现有定理的平凡推论——许多最终被纳入教科书的结果，同样在事前被同行预期过。因此“纲领非原创”只能否掉“新范式”这一说法，还不足以直接判定它是“凸闭包”。

\subsection{技术层：有一处实质性的新支点，但不是新范式}

判据二：跨越 \NS{} 这一具体障碍时，是否出现了既有纲领中没有、且不可由既有步骤机械合成的新构造步骤。

存在这样一处支点，也是本文认为最值得记录的地方：用零角平均振荡脉冲的 Reynolds 应力，去补足坍缩涡旋在奇点处缺失的动量输运。它的新意有两重。其一，它改变了被放大扰动的角色：在 Boussinesq 中，被放大的波本身就是奇点；在 \NS{} 中，波退居为工具，其平均自作用的散度充当作用在背景上的内部力。角色的重新指派不是尺度的微调，而是一次概念上的重定向。其二，它把“非线性项二次不消失”这一性质从障碍变成了资源：正是
\begin{equation}
\big\langle (w-\langle w\rangle)\otimes(w-\langle w\rangle)\big\rangle\neq 0
\end{equation}
这一事实，使得平均为零的扰动可以承担非零的通量。这项技术在更早的工作中已有先声（Daneri--Székelyhidi 用振荡实现指定应力、Craik--Criminale 的精确波、Córdoba--Martínez-Zoroa 的层次放大），但把“脉冲向坍缩核心补给动量”作为主线安排到 \NS{} 的强迫情形，是此处特有的组合。

与此配套的还有各向异性 Type II 坍缩：$\ell_{z}/\ell_{r}\asymp\tau^{-h}\to\infty$ 的细长柱，与以往层次构造的各向同性倾向不同。这一点的重要性后来被 CIV 的刚性定理反证式地确认——若核心不精确轴对称、或标度退化为各向同性，假想奇点将被逼回正则。因此这套形态不是任意的技术选择，而是构造得以成立的必要条件之一。

但需要同样清楚地指出：这些是技术层的支点，不是纲领层的突破；它们没有引入新的数学对象、新的不变量或新的普适结构，而是在既有纲领内部、为跨越一个具体障碍而设计的步骤。为把这一判断从断言变成可核验的结论，下面把构造的数学成分逐项拆开，核对其出处。

\subsection{成分溯源：五项可核验的出处}

\textbf{（一）迭代修正机制出自 Córdoba--Martínez-Zoroa。} 逐阶构造、把误差推向更高阶、让修正阶数随尺度上升，这些是 CMZ 纲领的规定动作。\NS{} 构造的第三至第十章是这套记账在更复杂对象上的执行，而非新机制。

\textbf{（二）精确波出自 Craik--Criminale，其不稳定性分析出自 Leibovich--Stewartson。} 第 4 章推出的振幅方程 $\dot a=-(I-2\,\xi\otimes\xi/|\xi|^{2})Da$ 与相位方程 $\dot\xi=-D^{\top}\xi$，在结构上与柱状涡的离心不稳定性分析完全一致：反射矩阵 $P_{\xi}=I-2\,\xi\otimes\xi/|\xi|^{2}$、振幅沿不稳定特征线的指数增长、波数随背景应变被拉伸，都是既有对象（见\cite{leibovich,billant}）。脉冲的生长没有引入新工具，只是把既有工具搬到环域上使用。

\textbf{（三）用振荡实现指定应力出自 Daneri--Székelyhidi，但用法不同。} 在凸积分传统中，振荡用于构造非唯一弱解、逼近一个 h-原理意义下的容许应力场，其目标落在“解的低正则自由度”上\cite{ds}。OpenAI 的用法在结构上不同：脉冲服务于一条\textbf{预先写定的光滑奇异背景}，其平均二次自作用的散度被用来抵消背景残差的奇异部分。这是“用同一材料、换一个工种”——振荡实现应力是既有的（棋类级）手段，而把应力供给绑定到一条坍缩背景上，是新的（棋步级）用法。这里随之出现一个此前没有的逆问题：给定必须实现的目标应力，构造能实现它的脉冲族；它的可解性由附录 C 的“容许应力锥”回答，而该附录在手段上仍属凸积分传统的技术（以振荡族的平均实现给定应力，可解性归结为某个锥条件），不引入新的数学对象。需要说明的是，本文未核对附录 C 的完整技术细节，这一归类以该附录的题名与相邻文献为据。

\textbf{（四）各向异性标度属于 CIV 已分析的 Type II 族，不是新标度类型。} $\ell_{r}\asymp\tau^{1/2}$、$\ell_{z}\asymp\tau^{1/2-h}$ 的奇点在轴向占优，比各向同性的 Type I（处处 $\tau^{-1/2}$）发散更快，因而落在 Type II 范畴；CIV 的刚性定理正是在这一族里工作，并证明只要再附加空间解析的外力，正则性就会回来\cite{civ}。更精确地说，CIV 在其注 1.4(b) 中明确指出，该定理与既有的一系列覆盖 Type I 情形的正则性结果——Chen--Strain--Yau--Tsai、Chen--Strain--Tsai--Yau、Koch--Nadirashvili--Seregin--Šverák、Seregin--Šverák 等——互不包含。这意味着各向异性 Type II 区间是先验定理的一块盲区：既有刚性结果既覆盖不到它、也排除不了它，而构造恰恰落在这块盲区里。因此这套标度是一个已被刻画过的族中的一点，其新颖性只在于“它是这条构造路径的必要选择”。

\textbf{（五）环面上的除子界与反演乘子是标准拟微分记账。} 引理 8.6 用到的两件工具——用除子界控制方向导数逆算子的象征、用二维 Fourier 可和性收尾——在拟微分算子与格点计数里都是常备件（第 8 章已给出完整推导）。它在这里的地位是“承重但普通”，不是新工具。

\textbf{屏障与转换：为何这一组合不可被机械合成。} 在五项溯源之外，还有一条来自方程本身的证据。CMZ 纲领的层次构造之所以能逐层升高频率，是因为在 Euler 与 IPM 中高频扰动没有黏性阻尼这个对手；而在 \NS{} 中，每个脉冲都背着 $-\nu|\xi|^{2}a$ 的衰减（第 6 章），频率越高衰减越快。这直接阻断了“把层放大原样译入 \NS{}”的机械路线，也是构造必须把扰动的角色从“奇点载体”改写为“应力供给者”的技术原因——角色转换不是风格选择，而是被黏性逼出来的。角色一旦改写，就必须回答那个此前没有的逆问题：给定必须实现的目标应力，构造能实现它的脉冲族（附录 C）。被黏性逼出的角色转换，加上随之而来的应力实现逆问题，合起来正是“不可机械合成”的具体含义。

还有一条时间线对照可以量化这道鸿沟。Albritton--Brué--Colombo\cite{abc}的非唯一性构造（2022）同样以“先选解、后定力”为骨架，其外力属于 $L_{t}^{1}L_{x}^{2}$ 且在初始时刻奇异；从这类允许奇异外力出场的结果，到备选 C 所要求的处处光滑、紧支、连同一切导数穿过奇点的外力，中间隔着的正是脉冲机制所要弥合的那段距离，也是这条路线此后数年主攻的方向。

把五项溯源与上述屏障证据并起来，得到一个比断言更强的判断：构造所用的每一块数学材料，都能在 2026 年之前已存在的文献里找到出处，没有任何一块是为此新造的数学对象、新不变量或新结构；而把这些既有材料拼装为有效证明的那一步——角色转换与应力实现——是可指认的具体新支点。这正是“它发明了一步新棋，而不是一种新棋类”的技术含义。

\subsection{方法论层：真正的新在流程，不在数学}

判据三：生产与验证知识的流程是否发生了阶跃。在这里，答案倾向于肯定，但必须明确这一肯定不落在数学思想层。

把一份由模型生成、据其共同作者自述“差到难以卒读”的写稿，在约一周内变成可被 Lean 内核编译的证书；把约一万个智能体组织成并行探索不同问题变体的搜索系统；用一次大模型调用完成对海量代码的形式化与验证——这些是若干年前不可想象的流程能力，它们改变的是求解与验证的成本结构。但按本层的判据，这属于方法论与工程，而不属于数学思想本身；把它记入“数学革命”的名下，会混淆两类不同的创新。

\subsection{综合判断}

把三层并起来，得到一个尽量不含混的定位。

在纲领层，它是既有纲领（Córdoba--Martínez-Zoroa 的残差博弈与逐尺度放大）在 \NS{} 强迫情形上的收官，思想内核非其原创。在技术层，它含有一处实质性的新支点（脉冲经 Reynolds 应力补缺力）与一套被刚性定理确认必要的各向异性坍缩形态，足以使它不落为机械合成，但尚不构成新范式。在方法论层，它是真正的新东西，也是本次事件中最可能被后人反复提及的部分，但它属于流程而非数学。

因此，若把“革命性的数学思想创新”定义为引入新对象、新结构或新不变量，并据此改写领域的问题框架，那么这项工作不是。若把它定义为“在既有纲领内，以一个此前没有的新支点跨过不可逾越的具体障碍”，那么它是，并且是其中规格相当高的一次。

\begin{table}[H]
\centering
\small
\caption{“思想定位”可操作判据表}
\begin{tabular}{@{}p{5.2cm}p{3.2cm}p{5.6cm}@{}}
\toprule
\hdrow \textbf{判据} & \textbf{本工作的答案} & \textbf{依据} \\
\midrule
是否引入新的数学对象或不变量 & 否 & 沿 Fefferman 既有提法与既有方程 \\
是否改变领域的问题框架 & 否 & 求解的是既有的备选 C、D \\
是否提出纲领之外的新原理 & 否 & Córdoba--Martínez-Zoroa 的残差博弈与逐尺度放大 \\
是否提供不可机械合成的关键步骤 & 是 & 脉冲经平均二次自作用补足坍缩缺力 \\
该步骤是否被后续工作确认为必要 & 是 & CIV 刚性定理排除各向同性与解析外力 \\
是否带来生产与验证流程的阶跃 & 是 & 多智能体搜索与规模化自动形式化（属方法论） \\
\bottomrule
\end{tabular}
\end{table}

最后需要一层诚实：上述判定以“我们确实理解了这份证明的动作”为前提，而这一点目前并不完全成立。形式化证书是自评的、尚未被独立核验；语义忠实性存在具体争议（第 8、9 章）；数学思想尚未被共同体严格重述。在证明被真正消化之前，“它在思想史中属于开创还是收官”这个问题只能给出暂定答案。这不是回避，而是对判据适用条件的如实交代。

\section{结语：修正后的判断}

本文开始时提出了一个问题：如何定位 OpenAI 在 \NS{} 上的这项工作。经过前十一章的梳理，可以给出一个尽量收敛、同时保留边界的判断。

\textbf{第一，它首先是一条技术路线的收敛证据。} 从 Córdoba--Martínez-Zoroa 的残差博弈，到 Alpöge--Buckmaster 的光滑外力与 Euler，再到把 \NS{} 的坍缩核心与由离心不稳定性生长的脉冲对接起来，这条路线有一以贯之的思想内核：先写下想要的奇异流，再让外力等于残差，然后集中全部技巧去保证残差光滑。OpenAI 的贡献在于把这套思想推到了 \NS{} 的强迫情形，其独有的一步是用零角平均脉冲的 Reynolds 应力补足坍缩缺失的动量输运。

\textbf{第二，它同时是一次对形式化验证认识的施压。} “Lean 通过”与“证明正确”之间，隔着“语义忠实翻译”这道在 SCI 谱系中高达无穷的坎；引理 8.6 处 $m+4$ 与 $m+5$ 的差异，是这道坎的一个具体样本。这条教训不针对 OpenAI，而适用于一切以自动形式化作为验证手段的工作。

\textbf{第三，它不是千禧年问题的终局裁判。} 备选 A 与 B 仍然开放；“有外力则可爆破”与“无外力是否永远光滑”本是两问。把它读成“AI 证明了水的方程会爆炸”，是对官方提法的误读。

\textbf{第四，边界必须写清楚。} 本文未独立复核该证明；关于正确性的判断均转述自公开文献并注明立场；CIV 的刚性定理是约束而非否证；争议事实均逐项标明来源、按归因分别叙述。这些限制不是谦辞，而是这篇综述得以成立的诚实前提。

把这四条并起来，得到的定位是：这是一次真实的、可检验的数学推进，也是一次对“什么才算验证”的追问。它推进的是一种构造方法，追问的是一种信念——我们是否愿意接受一个正确但尚未被理解的证明，作为数学知识被共同持有。

\begin{keybox}
\textbf{一句话收束}\quad 这项工作把一条酝酿多年的构造路线推到了 \NS{} 的强迫情形，也把“形式化通过是否等于理解发生”这个问题，第一次以如此具体的形态摊在数学共同体面前。前者是数学，后者是元数学，都还没有结束。
\end{keybox}

\appendix

\section{记号速查表}

\begin{table}[H]
\centering
\small
\caption{全文主要记号}
\begin{tabular}{@{}p{2.4cm}p{11.6cm}@{}}
\toprule
\hdrow \textbf{记号} & \textbf{含义} \\
\midrule
$u,p,f,\nu$ & 速度、压力、外力、黏性系数 \\
$\tau$ & 奇点前的剩余时间 $1-t$ \\
$\ell_{r},\ell_{z}$ & 核心的径向、轴向特征长度 \\
$h$ & 各向异性指数，$0<h<1/100$ \\
$\mathcal{P}u$ & 速度关于 $z$ 轴的角向平均（轴对称部分） \\
$w=u-\mathcal{P}u$ & 速度的非轴对称部分 \\
$\langle\cdot\rangle$ & 角向平均算子 \\
A/B/C/D & 千禧年问题官方提法的四条备选 \\
$\mathbb{T}^{3},\mathbb{T}^{2}$ & 三维、二维环面 \\
SCI & 可解性复杂度层级（Solvability Complexity Index） \\
\bottomrule
\end{tabular}
\end{table}

\section{一致性元数据}

\begin{table}[H]
\centering
\footnotesize
\caption{关键论断与可核验来源的对应}
\begin{tabular}{@{}p{5.2cm}p{8.8cm}@{}}
\toprule
\hdrow \textbf{本文论断} & \textbf{可核验来源} \\
\midrule
官方四备选 A--D 的表述 & Fefferman 撰写的 Clay 官方提法\cite{clay} \\
定理 1.1（全空间）与推论 10.6（周期） & OpenAI 手稿\cite{openai-paper} 第 1 节、第 10 节 \\
约一万并发智能体、88 小时、约 1300 亿输出 token & OpenAI 官方说明\cite{openai-blog} \\
Lean 仓库评审状态为“自评”、无 \texttt{sorry} & OpenAI Lean 仓库\cite{openai-lean}；Stanford Tech Review 审计\cite{stanford-audit} \\
Córdoba--Martínez-Zoroa 的 IPM 爆破 & \arxiv{2410.22920}\cite{cmz-ipm} \\
Alpöge--Buckmaster 的归属声明 & Buckmaster 公开声明\cite{buckmaster-statement} \\
各向异性标度 $\ell_{r}\asymp\tau^{1/2}$、$\ell_{z}\asymp\tau^{1/2-h}$ & OpenAI 手稿\cite{openai-paper} 第 2 节；CIV\cite{civ} 第 1 节 \\
引理 8.6 的 $m+4$ 与 Lean 的 $m+5$ & Bastounis--Circelli--Hansen\cite{bch} 例 3.1；并见本文据 OpenAI Lean 仓库\cite{openai-lean} commit \texttt{f9e8bc5} 的一手核验（第 8.3 节） \\
刚性定理（Type II ＋轴对称核心 ＋解析外力 $\Rightarrow$ 正则） & Constantin--Ignatova--Vicol\cite{civ} 定理 1.1 及其推论 \\
Clay 的“刻意不匆忙”表态 & Clay 新闻页\cite{clay-news} \\
正缺陷问题作为无外力靶标 & Petrillo--Glimm\arxiv{2609.23868}\cite{petrillo-glimm} \\
五个数学成分的出处（迭代修正、精确波、振荡实现应力、Type II 标度、环面乘子）与“屏障与转换”论证 & CMZ\cite{cmz-ipm}；Craik--Criminale 与 Leibovich--Stewartson\cite{leibovich,billant}；Daneri--Sz\'ekelyhidi\cite{ds}；CIV\cite{civ}；Albritton--Bru\'e--Colombo\cite{abc} \\
\bottomrule
\end{tabular}
\end{table}

\clearpage

\begin{thebibliography}{99}
\small
\setlength{\itemsep}{2pt}

\bibitem{clay} Fefferman C L. Existence and smoothness of the Navier--Stokes equation[C/OL]//The Millennium Prize Problems. Clay Mathematics Institute, 2000. \url{https://www.claymath.org/wp-content/uploads/2022/06/navierstokes.pdf}.

\bibitem{openai-blog} OpenAI. On the Navier--Stokes Millennium Prize Problem[EB/OL]. 2026-09-08. \url{https://openai.com/index/navier-stokes-solution/}.

\bibitem{openai-paper} OpenAI. Finite time blowup for Navier--Stokes[Z/OL]. 手稿, 2026. \url{https://cdn.openai.com/pdf/32d9f210-8b73-45e0-91bc-82a30aef8a9a/navier-stokes.pdf}.

\bibitem{openai-lean} OpenAI. NavierStokesAndEuler: Lean 4 formalizations[CP/OL]. GitHub. \url{https://github.com/openai/NavierStokesAndEuler}.

\bibitem{clay-news} Clay Mathematics Institute. Navier--Stokes announcement[EB/OL]. 2026-09-11. \url{https://www.claymath.org/news/navier-stokes-announcement/}.

\bibitem{leray} Leray J. Sur le mouvement d'un liquide visqueux emplissant l'espace[J]. Acta Mathematica, 1934, 63: 193--248.

\bibitem{ckn} Caffarelli L, Kohn R, Nirenberg L. Partial regularity of suitable weak solutions of the Navier--Stokes equations[J]. Communications on Pure and Applied Mathematics, 1982, 35(6): 771--831.

\bibitem{ess} Escauriaza L, Seregin G, \v{S}ver\'ak V. $L_{3,\infty}$-solutions of Navier--Stokes equations and backward uniqueness[J]. Russian Mathematical Surveys, 2003, 58(2): 211--250.

\bibitem{tao-averaged} Tao T. Finite time blowup for an averaged three-dimensional Navier--Stokes equation[J]. Journal of the American Mathematical Society, 2016, 29(3): 601--674.

\bibitem{bv-nonunique} Buckmaster T, Vicol V. Nonuniqueness of weak solutions to the Navier--Stokes equation[J]. Annals of Mathematics, 2019, 189(1): 101--144.

\bibitem{abc} Albritton D, Bru\'e E, Colombo M. Non-uniqueness of Leray solutions of the forced Navier--Stokes equations[J]. Annals of Mathematics, 2022, 196(1): 415--455.

\bibitem{ds} Daneri S, Sz\'ekelyhidi Jr L. Non-uniqueness and h-principle for H\"older-continuous weak solutions of the Euler equations[J]. Archive for Rational Mechanics and Analysis, 2017, 224(2): 471--514. DOI: 10.1007/s00205-017-1089-0.

\bibitem{cmz-ipm} C\'ordoba D, Mart\'inez-Zoroa L. Finite time singularities of smooth solutions for the 2D incompressible porous media (IPM) equation with a smooth source[Z/OL]. \arxiv{2410.22920}, 2024.

\bibitem{cmz-euler} C\'ordoba D, Mart\'inez-Zoroa L, Zheng F. 光滑外力下流体方程奇点形成的相关工作[Z]. 见手稿\cite{openai-paper} 参考文献 [6][8]。

\bibitem{buckmaster-statement} Buckmaster T. Statement[EB/OL]. 2026-09-08. \url{https://cims.nyu.edu/~tristanb/statement.pdf}.

\bibitem{tao-blog-sep7} Tao T. Finite time blowup with smooth forcing term for the incompressible porous medium, Boussinesq, and incompressible Euler equations[EB/OL]. What's new, 2026-09-07. \url{https://terrytao.wordpress.com/2026/09/07/}.

\bibitem{schneider} Schneider T. Headlines and inside stories: understanding and trust in AI for mathematics, science, and engineering[EB/OL]. 客座于 What's new, 2026-09-23. \url{https://terrytao.wordpress.com/2026/09/23/}.

\bibitem{gowers} Gowers T. Why I didn't sign the Fields medallists' letter[EB/OL]. 2026-09-17. \url{https://terrytao.wordpress.com/2026/09/17/}.

\bibitem{mathandai} 25 Fields medallists. Letter on the misalignment between AI companies and the mathematical community[EB/OL]. \url{https://mathandai.org/}.

\bibitem{civ} Constantin P, Ignatova M, Vicol V. Regularity of asymptotically axisymmetric solutions to the 3D Navier--Stokes equations with analytic forcing[Z/OL]. \arxiv{2609.20803}, 2026.

\bibitem{bch} Bastounis A, Circelli F, Hansen A C. Navier--Stokes lost in translation: why Lean verification of AI autoformalisation does not guarantee correct natural language proofs[Z/OL]. \arxiv{2610.08144}, 2026.

\bibitem{petrillo-glimm} Petrillo J, Glimm J. The positive defect problem: target and admissibility criteria for a programmatic search for unforced Navier--Stokes blowup[Z/OL]. \arxiv{2609.23868}, 2026.

\bibitem{shavit} Shavit M. Two forced finite-time singularities, and the mechanism behind each: a friendly reading of Alp\"oge--Buckmaster (Boussinesq) and OpenAI (Navier--Stokes)[EB/OL]. 2026-09. \url{https://michalshavitnyu.github.io/files/blowup-mechanisms.pdf}.

\bibitem{stanford-audit} Raman P. OpenAI vs Buckmaster: the Navier--Stokes Lean proofs, audited[EB/OL]. Stanford Tech Review, 2026-09-08. \url{https://stanfordtechreview.com/articles/openai-buckmaster-navier-stokes-lean-proofs}.

\bibitem{leibovich} Leibovich S, Stewartson K. A sufficient condition for the instability of columnar vortices[J]. Journal of Fluid Mechanics, 1983, 126: 335--356.

\bibitem{billant} Billant P, Gallaire F. Generalized Rayleigh criterion for non-axisymmetric centrifugal instabilities[J]. Journal of Fluid Mechanics, 2005, 542: 365--379.

\end{thebibliography}

\end{document}
